\chapter{知识谱系图示}\label{ch:spectrum_map}
\label{ch:spectrum_diagram}

% ============================================================
\section*{本章导读}
\addcontentsline{toc}{section}{本章导读}
% ============================================================

本章\textbf{不引入新定理}。
第~\ref{ch:riemann_explicit_formula}~章完成显式公式与 $M(x)$ 等价后，
主体理论链条（初等算术 $\to$ $\zeta$ 延拓 $\to$ $\xi$ 与零点 $\to$ 显式公式）已齐备。

全书共用\textbf{同一底图}，分\textbf{三次着色}标示阅读进度
（彩色 $=$ 截至该进度已建立的主要对象与关系，灰白 $=$ 后文）：
\begin{enumerate}
  \item \textbf{第一次}：第~\ref{ch:elementary_number_theoretic_functions}~章末，
    初等算术层（图~\ref{fig:spectrum-stage2}）；
  \item \textbf{第二次}：第~\ref{ch:riemann_zeta_continuation}~章末，
    叠加 $\Gamma$、$\zeta$ 与延拓/函数方程（图~\ref{fig:spectrum-stage8}）；
  \item \textbf{第三次}（本章）：全文着色总图（图~\ref{fig:function-relations}）。
\end{enumerate}
本章为第三次着色，作为进入数值、几何示意与原稿解读之前的导航索引。
图上箭头应读作已证公式或定义的索引，\textbf{不替代}各章证明。

% ============================================================
\section{全文着色总图}
\label{sec:spectrum_full_map}
% ============================================================

至此，A-2 底图上与主体理论相关的对象均已在正文中建立：
第~\ref{ch:elementary_number_theoretic_functions}~章的初等算术层
（图~\ref{fig:spectrum-stage2}）、
第~\ref{ch:riemann_zeta_continuation}~章的 $\Gamma$、$\zeta$ 与延拓/函数方程
（图~\ref{fig:spectrum-stage8}），以及第~\ref{ch:xi_function}~、
\ref{ch:zero_distribution}~章的 $\xi$、$N(T)$、RH 表述，
与第~\ref{ch:riemann_explicit_formula}~章的 Perron/Mellin 显式公式及 $M(x)$ 等价。
本章将底图上的\textbf{主要对象与关系一次着色}，完成上述三次进度中的最后一次。
（图为选录：只标主体理论脉络中的主要节点与连线，并非穷尽一切对象与关系。）
第一次、第二次图元索引见表~\ref{tab:spectrum-stage2-guide}、
表~\ref{tab:spectrum-stage8-guide}；下表仅列\textbf{相对第二次进度图（图~\ref{fig:spectrum-stage8}）新着色}者。

% 图独占一页；题注文字排到下一页，避免与图例挤在同页
\clearpage
\begin{figure}[p]
\centering
\def\spectrumStage{16}
\resizebox{\textwidth}{!}{%
% 黎曼猜想知识谱系关系图（A-2 现行底图 + 分阶段灰白）
% 用法：\def\spectrumStage{2|8|16}% 黎曼猜想知识谱系关系图（A-2 现行底图 + 分阶段灰白）
% 用法：\def\spectrumStage{2|8|16}% 黎曼猜想知识谱系关系图（A-2 现行底图 + 分阶段灰白）
% 用法：\def\spectrumStage{2|8|16}\input{图谱/谱系关系图-core.tex}
% 几何与 黎曼猜想知识谱系关系图-A-2.tex 同步；未涉及者仅灰白化（含图例）
%
% \spectrumStage 仅为内部进度码（历史命名），不是章号：
%   2  → 第一次着色：第2章 初等算术层（fig:spectrum-stage2）
%   8  → 第二次着色：第9章 延拓/函数方程（fig:spectrum-stage8）
%  16  → 第三次着色：第15章 全文总图（fig:function-relations）
% 正文表述用「第一次 / 第二次 / 第三次着色」或章标签，勿写「第8章图」「第16章图」。
\providecommand{\spectrumStage}{16}
\makeatletter
\colorlet{cDim}{gray!16}\colorlet{dDim}{gray!48}\colorlet{tDim}{gray!52}
\ifnum\spectrumStage=2
  \colorlet{cPrime}{blue!12}\colorlet{dPrime}{black}\colorlet{tPrime}{black}
  \colorlet{cMu}{teal!15}\colorlet{dMu}{teal!60!black}\colorlet{tMu}{black}
  \colorlet{cLi}{yellow!25}\colorlet{dLi}{orange!70!black}\colorlet{tLi}{black}
  \colorlet{cGammaU}{gray!16}\colorlet{cGammaL}{gray!16}\colorlet{dGamma}{gray!48}\colorlet{tGamma}{gray!52}
  \colorlet{cZetaU}{gray!16}\colorlet{cZetaL}{gray!16}\colorlet{tZeta}{gray!52}
  \colorlet{cXi}{gray!16}\colorlet{dXi}{gray!48}\colorlet{tXi}{gray!52}
  \colorlet{cNT}{gray!16}\colorlet{dNT}{gray!48}\colorlet{tNT}{gray!52}
  \colorlet{cProd}{gray!16}\colorlet{dProd}{gray!48}\colorlet{tProd}{gray!52}
  \colorlet{cRH}{gray!16}\colorlet{dRH}{gray!48}\colorlet{tRH}{gray!52}
  \colorlet{eArith}{blue!70!black}\colorlet{eElem}{violet!80!black}\colorlet{eDef}{black!80}
  \colorlet{ePNT}{brown!70!black}\colorlet{eAnal}{gray!42}\colorlet{eExpl}{gray!42}
  \colorlet{eExt}{gray!42}\colorlet{eZero}{gray!42}\colorlet{eLate}{gray!42}
  \colorlet{lgPrime}{blue!12}\colorlet{lgMu}{teal!15}\colorlet{lgLi}{yellow!25}
  \colorlet{lgXi}{gray!22}\colorlet{lgRH}{gray!22}\colorlet{lgGU}{gray!22}\colorlet{lgGL}{gray!22}
  \colorlet{lgEArith}{blue!70!black}\colorlet{lgEDef}{black!80}\colorlet{lgEElem}{violet!80!black}
  \colorlet{lgEZero}{gray!45}\colorlet{lgEAnal}{gray!45}\colorlet{lgEPNT}{brown!70!black}\colorlet{lgEExt}{gray!45}
  \colorlet{lgTprime}{black}\colorlet{lgTmu}{black}\colorlet{lgTli}{black}
  \colorlet{lgTxi}{gray!50}\colorlet{lgTrh}{gray!50}\colorlet{lgTgamma}{gray!50}
  \colorlet{lgTarith}{black}\colorlet{lgTdef}{black}\colorlet{lgTelem}{black}
  \colorlet{lgTzero}{gray!50}\colorlet{lgTanal}{gray!50}\colorlet{lgTpnt}{black}\colorlet{lgText}{gray!50}
\else
\ifnum\spectrumStage=8
  \colorlet{cPrime}{blue!12}\colorlet{dPrime}{black}\colorlet{tPrime}{black}
  \colorlet{cMu}{teal!15}\colorlet{dMu}{teal!60!black}\colorlet{tMu}{black}
  \colorlet{cLi}{yellow!25}\colorlet{dLi}{orange!70!black}\colorlet{tLi}{black}
  \colorlet{cGammaU}{blue!15}\colorlet{cGammaL}{red!15}\colorlet{dGamma}{black}\colorlet{tGamma}{black}
  \colorlet{cZetaU}{blue!15}\colorlet{cZetaL}{red!15}\colorlet{tZeta}{black}
  \colorlet{cXi}{gray!16}\colorlet{dXi}{gray!48}\colorlet{tXi}{gray!52}
  \colorlet{cNT}{gray!16}\colorlet{dNT}{gray!48}\colorlet{tNT}{gray!52}
  \colorlet{cProd}{gray!16}\colorlet{dProd}{gray!48}\colorlet{tProd}{gray!52}
  \colorlet{cRH}{gray!16}\colorlet{dRH}{gray!48}\colorlet{tRH}{gray!52}
  \colorlet{eArith}{blue!70!black}\colorlet{eElem}{violet!80!black}\colorlet{eDef}{black!80}
  \colorlet{ePNT}{brown!70!black}\colorlet{eAnal}{green!50!black}\colorlet{eExpl}{gray!42}
  \colorlet{eExt}{red!70!black}\colorlet{eZero}{gray!42}\colorlet{eLate}{gray!42}
  \colorlet{lgPrime}{blue!12}\colorlet{lgMu}{teal!15}\colorlet{lgLi}{yellow!25}
  \colorlet{lgXi}{gray!22}\colorlet{lgRH}{gray!22}\colorlet{lgGU}{blue!15}\colorlet{lgGL}{red!15}
  \colorlet{lgEArith}{blue!70!black}\colorlet{lgEDef}{black!80}\colorlet{lgEElem}{violet!80!black}
  \colorlet{lgEZero}{gray!45}\colorlet{lgEAnal}{green!50!black}\colorlet{lgEPNT}{brown!70!black}\colorlet{lgEExt}{red!70!black}
  \colorlet{lgTprime}{black}\colorlet{lgTmu}{black}\colorlet{lgTli}{black}
  \colorlet{lgTxi}{gray!50}\colorlet{lgTrh}{gray!50}\colorlet{lgTgamma}{black}
  \colorlet{lgTarith}{black}\colorlet{lgTdef}{black}\colorlet{lgTelem}{black}
  \colorlet{lgTzero}{gray!50}\colorlet{lgTanal}{black}\colorlet{lgTpnt}{black}\colorlet{lgText}{black}
\else
  \colorlet{cPrime}{blue!12}\colorlet{dPrime}{black}\colorlet{tPrime}{black}
  \colorlet{cMu}{teal!15}\colorlet{dMu}{teal!60!black}\colorlet{tMu}{black}
  \colorlet{cLi}{yellow!25}\colorlet{dLi}{orange!70!black}\colorlet{tLi}{black}
  \colorlet{cGammaU}{blue!15}\colorlet{cGammaL}{red!15}\colorlet{dGamma}{black}\colorlet{tGamma}{black}
  \colorlet{cZetaU}{blue!15}\colorlet{cZetaL}{red!15}\colorlet{tZeta}{black}
  \colorlet{cXi}{cyan!15}\colorlet{dXi}{cyan!60!black}\colorlet{tXi}{black}
  \colorlet{cNT}{cyan!15}\colorlet{dNT}{cyan!60!black}\colorlet{tNT}{black}
  \colorlet{cProd}{cyan!15}\colorlet{dProd}{cyan!60!black}\colorlet{tProd}{black}
  \colorlet{cRH}{orange!15}\colorlet{dRH}{orange!70!black}\colorlet{tRH}{black}
  \colorlet{eArith}{blue!70!black}\colorlet{eElem}{violet!80!black}\colorlet{eDef}{black!80}
  \colorlet{ePNT}{brown!70!black}\colorlet{eAnal}{green!50!black}\colorlet{eExpl}{green!50!black}
  \colorlet{eExt}{red!70!black}\colorlet{eZero}{orange!80!black}\colorlet{eLate}{black!80}
  \colorlet{lgPrime}{blue!12}\colorlet{lgMu}{teal!15}\colorlet{lgLi}{yellow!25}
  \colorlet{lgXi}{cyan!15}\colorlet{lgRH}{orange!15}\colorlet{lgGU}{blue!15}\colorlet{lgGL}{red!15}
  \colorlet{lgEArith}{blue!70!black}\colorlet{lgEDef}{black!80}\colorlet{lgEElem}{violet!80!black}
  \colorlet{lgEZero}{orange!80!black}\colorlet{lgEAnal}{green!50!black}\colorlet{lgEPNT}{brown!70!black}\colorlet{lgEExt}{red!70!black}
  \colorlet{lgTprime}{black}\colorlet{lgTmu}{black}\colorlet{lgTli}{black}
  \colorlet{lgTxi}{black}\colorlet{lgTrh}{black}\colorlet{lgTgamma}{black}
  \colorlet{lgTarith}{black}\colorlet{lgTdef}{black}\colorlet{lgTelem}{black}
  \colorlet{lgTzero}{black}\colorlet{lgTanal}{black}\colorlet{lgTpnt}{black}\colorlet{lgText}{black}
\fi\fi
\makeatother

\begin{tikzpicture}[scale=1.0, transform shape,
  box/.style    = {draw, rounded corners=5pt, minimum width=2.8cm,
                   minimum height=0.9cm, align=center, font=\small, thick},
  zbox/.style   = {draw, rounded corners=5pt, minimum width=2.8cm,
                   minimum height=0.9cm, align=center, font=\small,
                   thick, fill=gray!18},
  splitbox/.style = {draw, rounded corners=5pt, minimum width=7cm,
                     minimum height=2.2cm, align=center, font=\small, thick,
                     rectangle split, rectangle split parts=2,
                     rectangle split part fill={red!8, white}},
  arr/.style    = {-{Stealth[length=6pt,width=5pt]}, thick},
  barr/.style   = {{Stealth[length=6pt,width=5pt]}-{Stealth[length=6pt,width=5pt]}, thick},
  % 线型：虚线仅用于 N(T)→RH、ξ↔RH；其余为实线
  proven/.style = {solid},
  cond/.style   = {dashed},
  lab/.style    = {font=\tiny, align=center, inner sep=3pt},
  rhbox/.style  = {draw=dRH, rounded corners=5pt, minimum width=8cm,
                   minimum height=1.0cm, align=center, font=\small,
                   fill=cRH, text=tRH, thick}
]

% 布局与 A-2 同步：左移 1.0cm、上移 2.0cm
% 左右略放宽，避免 Γ/ζ 自环旁公式标注被裁切
\path[use as bounding box]
  (-10.2, -26.2) rectangle (10.0, 7.4);

\begin{scope}[shift={(-1.0,2.0)}]

\node[font=\bfseries\Large, align=center] at (0, 4.35)
  {黎曼猜想知识谱系关系图};

% 主图单独上移 1cm；标题与图例位置不动
\begin{scope}[shift={(-2.75,-1.2)}]
%% ==================== 节点分层 ====================

% 第一层：素数计数函数 (y=0)
\node[box, fill=cPrime, draw=dPrime]    (J)    at (-4.2,  0.0) {$J(x)$};
\node[box, fill=cPrime, draw=dPrime]    (pi)   at ( 0.0,  0.0) {$\pi(x)$};
\node[box, fill=cPrime, draw=dPrime]    (th)   at ( 5.0,  0.0) {$\theta(x)$};
\node[box, fill=cPrime, draw=dPrime]    (ps)   at ( 9.7,  0.0) {$\psi(x)$};

% 第二层：积性函数 (y=-3.0)
% 节点定义修改
\node[box, fill=cMu, draw=dMu] (mu)  at (0.3, -3.0) {$\mu(n)$};
\node[box, fill=cMu, draw=dMu] (Lam) at (5.0, -3.0) {$\Lambda(n)$};

%Weierstrass乘积
\node[box, fill=cProd, draw=dProd, text=tProd]  (prod)   at ( 9.5,  -17.0) {\tiny\text{Hadamard 乘积}};
% M(x) 节点已移除：初等等价信息直接标注在 RH 框内

% 第三层：Li 与 Gamma (y=-6.0 至 -7.0)
\node[box, fill=cLi, draw=dLi] (Li) at (-2.0, -6.0) {$\Li(x)$};

% Gamma(s) 上下分色：上部标收敛域，下部标延拓
\node[splitbox, rectangle split draw splits=false, minimum width=5cm,
      minimum height=4.7cm,
      rectangle split part fill={cGammaU, cGammaL}] (gamma) at (7.0, -6.0)
    {
     \centering \color{tGamma}$\operatorname{Re}(s)>0$
     \nodepart{two}
     \raisebox{6pt}{\color{tGamma}$\mathbb{C}\setminus\{0,-1,-2,\dots\}$}
    };
% Gamma(s) 标注（固定在节点中线左侧）
\node[font=\small, left=-30pt, text=tGamma] at (gamma.west) {$\Gamma(s)$};

% 【修订3】ζ(s) 节点：下半部分补全函数方程左侧 ζ(s)=
\node[splitbox, rectangle split draw splits=false, minimum height=4cm,
      rectangle split part fill={cZetaU, cZetaL}] (zeta) at (0.3, -10.5)
    {\centering\small
     \color{tZeta}$\displaystyle
       \sum_{n=1}^{\infty}\frac{1}{n^s}
       =\prod_{p}\dfrac{1}{1-p^{-s}}
       \quad(\operatorname{Re}(s)>1)$%
     \vphantom{$\displaystyle\sum_{n=1}^{\infty}$}%
     \nodepart{two}
     % 修订：补全 ζ(s)= 使函数方程完整
     \textcolor{tZeta}{$\displaystyle\zeta(s)=
       2^s\pi^{s-1}
       \sin\!\left(\frac{\pi s}{2}\right)\Gamma(1-s)\,\zeta(1-s)$%
     \vphantom{$\displaystyle\sum_{n=1}^{\infty}\mathbb{C}\setminus\{1\}$}}
    };
% ζ(s) 标注


% ξ(s) 节点：下移至与 RH 同一水平，右侧对称于 M(x)
\node[box, fill=cXi, draw=dXi, text=tXi, minimum height=1.2cm]
  (xi) at (4.85, -17.0)
  {$\xi(s)$\\ \tiny 非平凡零点$\rho$与$\zeta(s)$相同\\
  \tiny $\xi(s)=\xi(1-s)$(函数方程推论)
  };


% 【修订8】N(T) 节点：置于 ζ 与 RH 正中间，构成竖直逻辑链
\node[box, fill=cNT, draw=dNT, text=tNT, minimum height=1.3cm]
  (NT) at (2.3, -14.0)
  {$N(T)$\\ \tiny $\zeta(s)$ 零点计数函数\\[1pt]\tiny（幅角原理）};

% ==================== 黎曼猜想节点 ====================
% RH 下移至 y=-20.5，为 N(T) 留出充分空间
\node[rhbox, minimum width=4.6cm, minimum height=2.2cm] (rh) at (-1.9, -17.0) {
    {\tiny\color{tRH} 黎曼猜想 (Riemann Hypothesis, 1859)}\\[3pt]
    {\tiny\color{tRH}$\zeta(s)$ 的所有非平凡零点 $\rho$ 满足
    $\operatorname{Re}(\rho)=\dfrac{1}{2}$}\\[4pt]
{\tiny\color{tRH} 初等等价（梅滕斯函数）：%
  $\displaystyle M(x)=\sum_{n\leq x}\mu(n)$}\\
{\tiny\color{tRH}
  $\mathrm{RH} \iff M(x) = O(x^{1/2+\varepsilon}),\ \forall \varepsilon>0$}
};


% ==================== 连线 ====================

% ----- 第一层内部 -----
\draw[barr, eArith] (J.east) -- (pi.west)
  node[lab, midway, above=14pt, xshift=-2pt, eArith]
    {$\displaystyle\pi(x)=\sum_{n=1}^{\infty}\frac{\mu(n)}{n}J(x^{1/n})$}
  node[lab, midway, below=14pt, xshift=-2pt, eArith]
    {$\displaystyle J(x)=\sum_{n=1}^{\infty}\frac{1}{n}\pi(x^{1/n})$};

\draw[barr, eElem] (pi.east) -- (th.west)
  node[lab, midway, above=14pt, eElem]
    {$\displaystyle\theta(x)=\pi(x)\log x-\int_2^x\frac{\pi(t)}{t}\,\mathrm{d}t$}
  node[lab, midway, above={14pt}, yshift={-50pt}, eElem]
    {\tiny $\displaystyle\pi(x) = \dfrac{\theta(x)}{\log x}+ \int_{2}^{x} \frac{\theta(t)}{t\,(\log t)^2} \,\mathrm{d}t$};

\draw[barr, eElem] (th.east) -- (ps.west)
  node[lab, midway, above=11pt, xshift={5pt}, eElem]
    {$\displaystyle\psi(x)=\sum_{k=1}^{\infty}\theta(x^{1/k})$}
  node[lab, midway, above=-31pt, xshift={5pt}, yshift={-8pt}, eElem]
    {\tiny $\displaystyle\theta(x)=\sum_{k=1}^{\infty}\mu(k)\psi(x^{1/k})$};


% ----- 第二层内部 -----
\draw[barr, eArith] (mu.east) -- (Lam.west)
  node[lab, midway, yshift=20, xshift=0.3cm, eArith]
    {$\displaystyle\Lambda(n)=-\sum_{d\mid n}\mu(d)\log d$
     \quad{\rm (M\"{o}bius 反演)}}
  node[lab, midway, yshift={-12pt-0.5cm}, xshift=0, eArith]
    {$\displaystyle\log n=\sum_{d\mid n}\Lambda(d)$};

% μ(n) → M(x) 的定义连线已省去（见节点内文字），避免与已有连线交叉

% Lambda -> psi
\draw[arr, eDef] (Lam.east) to[out=0, in=-90, looseness=1.0]
  ($(ps.south west)+(30pt,0)$)
  node[lab, pos=0.5, yshift=-90, xshift=250, eDef]
    {$\displaystyle\psi(x)=\sum_{n\le x}\Lambda(n)$};

% ----- 第三层：Li 和 Gamma -----

% 【修订5】Li(x) → π(x)：改为渐近等价语义（非因果推导）
\draw[arr, ePNT]
  ($(Li.north)!.75!(Li.north east)$)
  to[out=90, in=-135, looseness=1.8]
  ($(pi.south west)+(35pt,0)$)
  node[lab, pos=0.4, yshift=-45, xshift=45, ePNT, align=left]
    {\textbf{素数定理（PNT）}：$\pi(x)\sim\Li(x)$\\
     (若RH成立，$\pi(x)=\Li(x)+O(\sqrt{x}\log x)$)};

% Li -> J（显式公式主项；属 eExpl：第2/9章灰白，第15章全文总图着色）
\draw[arr, eExpl]
  (Li.north) to[out=135, in=-90, looseness=1.5]
  ($(J.south east)+(-22pt,0)$)
  node[lab, pos=0.5, yshift=-115, xshift=-65, eExpl, align=center]
    {$\Li(x)$ 是 $J(x)$ 显式\\[2pt]公式的主项};

% Gamma（Re>0）→ ζ：Euler 积分表示（该积分公式本身要求 Re(s)>1）
\draw[arr, eAnal]
  ([yshift=6pt] gamma.west)
  to[out=180, in=90, looseness=0.8]
  ($(zeta.north)+(30pt,0)$)
  node[lab, pos=0.5, yshift=-235, xshift=90,eAnal, align=center]
    {$\displaystyle\zeta(s)\Gamma(s)
      =\int_0^{\infty}\frac{x^{s-1}}{e^x-1}\,\mathrm{d}x$\\
     $(\operatorname{Re}(s)>1)$};

% μ(n) → ζ(s)
\draw[arr, eAnal]
  (mu.south) to[out=-90, in=90, looseness=1.2]
  ($(zeta.north)+({-20pt+0.3cm},0)$)
  node[lab, pos=0.5, yshift={-120pt-0.5cm-0.4cm}, xshift={45pt-5pt}, eAnal]
    {$\displaystyle\frac{1}{\zeta(s)}=\sum_{n=1}^{\infty}\frac{\mu(n)}{n^s}$\\
     $(\operatorname{Re}(s)>1)$};

% ----- J(x) 与 ψ(x) 的 Stieltjes 积分互逆关系 -----
% 【修订4】加注 dψ(t)/ln t 在 t=1 附近需 t>1 的说明
\draw[barr, eArith, line width=0.8pt, rounded corners=8pt]
  ($(J.north)+(0,0.0cm)$)
  -- ++(0,2.0cm)
  -- ++(13.5cm,0)
  -- ++(0,-2.0cm)
  node[lab, pos=0.35, above=6pt, yshift={10pt+0.5cm}, xshift=-200, eArith, align=center]
    {$\displaystyle\psi(x)=\int_{2^-}^x\log t\,\mathrm{d}J(t)$\quad(Stieltjes 积分，$J\to\psi$)}
  node[lab, pos=0.65, below=6pt, yshift={32pt+0.2cm}, xshift=-205, eArith, align=center]
    {$\displaystyle J(x)=\int_{2^-}^x\frac{\mathrm{d}\psi(t)}{\log t}$
     \quad(Stieltjes 反演，$\psi\to J$；积分自 $t>1$)};

% ----- 第四层：ζ 和 ξ 与第二层的连线 -----

% ζ → ψ：显式公式（Perron 积分 + 留数，单向；x≠p^k 时公式精确成立）
\draw[arr, eExpl, rounded corners=30pt]
  ($(zeta.east)+(0, -10pt)$) -| ($(ps.south)+(30pt,0)$)
  node[lab, pos=0.85, right=10pt, yshift=-170, xshift={-170pt-0.2cm},eExpl, align=left]
{$\psi(x)$ 显式公式（Perron 积分，$\zeta\to\psi$）\\
     $\displaystyle \psi(x) = \frac{1}{2\pi i}\int_{c-i\infty}^{c+i\infty}\!\left(-\frac{\zeta'(s)}{\zeta(s)}\right)\!\frac{x^s}{s}\,\mathrm{d}s$\\
     $\displaystyle = x-\sum_{\rho}\frac{x^{\rho}}{\rho}-\log(2\pi)-\tfrac{1}{2}\log(1-x^{-2})$\\
     \tiny $x>1,\;x\neq p^k$；$\rho$ 按 $|\mathrm{Im}(\rho)|$ 升序配对，条件收敛};

% ψ → ζ：Mellin 变换（单向）
\draw[arr, eExpl, rounded corners=30pt]
  ($(ps.south)+(40pt,0)$) |- ($(zeta.east)+(0, -20pt)$)
  node[lab, pos=0.8, yshift=-27, xshift=30,eExpl, align=center]
    {Mellin 变换（$\psi(x) \longleftrightarrow -\dfrac{\zeta'(s)}{\zeta(s)}$，即负对数导数）\\
     $\displaystyle -\frac{\zeta'(s)}{\zeta(s)} = s\int_{1}^{\infty}\frac{\psi(x)}{x^{s+1}}\,\mathrm{d}x$};
       
       
       
       
       
       
       
% Lambda -> zeta（对数微分 Dirichlet 级数）
\draw[arr, eAnal]
  ($(Lam.south west)+(6pt,0)$) to[out=-90, in=90, looseness=0.4]
  ($(zeta.north west)!0.50!(zeta.north east)$)
  node[lab, pos=0.5, yshift={-120pt-0.5cm-0.4cm+1cm}, xshift={145pt-5pt}, eAnal]
    {$\displaystyle-\frac{\zeta'(s)}{\zeta(s)}
      =\sum_{n=1}^{\infty}\frac{\Lambda(n)}{n^s}$};

% J(x) 显式公式（ζ → J）
\draw[arr, eExpl, rounded corners=50pt]
  ($(zeta.south west)+(0,6pt)$) -- ++(-60pt,0) -- ($(J.south west)+(5pt,0)$)
  node[lab, pos=0.5, yshift=-203, xshift={75pt-0.2cm}, eExpl, align=center]
{\textbf{Perron 公式} $\Rightarrow$ $J(x)$ 显式展开\\
     $\displaystyle
     \begin{aligned}
         &J(x)=\frac{1}{2\pi i}\int_{c-i\infty}^{c+i\infty}
           \frac{\log\zeta(s)}{s}\,x^s\,\mathrm{d}s\\
         &=\Li(x)-\sum_{\rho}\Li(x^{\rho})
           -\log 2+\int_x^{\infty}\frac{\mathrm{d}t}{t(t^2-1)\log t}
     \end{aligned}$\\
     \tiny $x>1,\;x\neq p^k$；$\rho$ 按 $|\mathrm{Im}(\rho)|$ 升序配对，条件收敛};

% ζ → J：Mellin 变换（ln ζ）
\draw[arr, eAnal, rounded corners=5pt]
  (zeta.north west) to[out=110, in=-90, looseness=1.5] (J.south)
  node[lab, pos=0.4, yshift=-205, xshift=-55, eAnal, align=center]
    {$\displaystyle\log\zeta(s)=s\int_1^{\infty}J(x)\,x^{-s-1}\,\mathrm{d}x$\\
     $(\operatorname{Re}(s)>1)$};

% ζ → ξ（定义）：ξ 已下移，连线沿右侧空白区域向右下延伸
\draw[arr, eLate, rounded corners=15pt]
  ($(zeta.south east)!0.15!(zeta.north east)$)
  -- ++(1.0, 0)
  -- ++(0, -3.8)
  -- ($(xi.north)+(0pt,0)$)
  node[lab, pos=0.25, right=4pt,yshift=30, xshift=0, eLate, align=center]
    {$\xi(s)=\dfrac{1}{2}s(s-1)\pi^{-s/2}\Gamma(s/2)\zeta(s)$  
};






% xi -> prod (Weierstrass product connection)
\draw[arr, eLate]
  %(xi.south) -- (prod.south)
  (xi.east) --(prod.west)
  node[midway, right=4pt, yshift=28, xshift=-45, eLate]
    {\tiny$\displaystyle\xi(s)=e^{A+Bs}\prod_{\rho}\left(1-\frac{s}{\rho}\right)e^{s/\rho}$};

% prod -> ps (zero sum in explicit formula)
%\draw[arr, violet!70!black, rounded corners=15pt]
  %(prod.east) -- ++(4.4, 0) |- (ps.south)
  %node[lab, pos=0.6, right=2pt, violet!70!black, align=left]
   % {对数微分与 Perron 积分\\
    % 给出显式公式中\\
    % 零点和 $\sum_{\rho}\frac{x_{\rho}}{\rho}$ 项};
%
% ζ 自身左侧弧线（解析延拓）：与表 tab:spectrum-stage8-guide 一致
\draw[arr, eExt, rounded corners=5pt]
  ($(zeta.north west)!0.15!(zeta.south west)$)
  to[out=180, in=180, looseness=1.6]
  ($(zeta.north west)!0.85!(zeta.south west)$);
% 标签用独立节点锚定在 ζ 西侧弧线旁（避免 path 末 node 漂移）
% 单行；中心在 zeta.west 右 2.5cm（相对初版再累计右移）
\node[lab, eExt, align=center, anchor=center]
  at ([xshift=2.5cm]zeta.west)
  {$\zeta$ 亚纯延拓至 $\mathbb{C}\setminus\{1\}$（Hankel / $\vartheta$ 两法）};

% ζ → N(T)：竖直向下（弧线）
\draw[arr, eZero]
  ($(zeta.south east)!0.10!(zeta.north east)$)
  to[out=0, in=90] ++(0.8, -0.8)
  to[out=-90, in=0] (NT.east)
  node[lab, pos=0.6, right=6pt, yshift=-357, xshift=5, eZero, align=center]
    {$\displaystyle N(T)=\frac{T}{2\pi}\log\frac{T}{2\pi e}+O(\log T)$\\
     \tiny 零点计数公式（von Mangoldt）};

% 【已删去】ζ → RH 蓝色直连线（设计意图改为ζ→N(T)→RH线性链）
%\draw[arr, eArith, rounded corners=10pt]
%  (zeta.south)
%  -- ++(0, -3.35)
%  -- ++(-4.1, 0)
%  -- ($(rh.north)!0.80!(rh.north west)$)
%  node[lab, pos=0.35, left=6pt, yshift=20pt,xshift=120,eArith, align=center]
%    {
%     $\zeta(s)$ 的非平凡零点全部位于直线 $\operatorname{Re}(s)=\dfrac12$ 上};
% RH -- P(s) 画线 + 文字
%\draw[arr, cyan!60!black]
 % ($(rh.east)!0.60!(rh.south east)$) -- (prod.west)    % 起点 -- 终点（必须有！）
  %node[midway, right=6pt, font=\tiny, cyan!60!black, align=center]
  %  {$\zeta(s) = \frac{e^{s(\log(2\pi) - 1)}}{2(s-1)} \prod_{n=1}^{\infty} \left(1 + \frac{s}{2n}\right)e^{-s/2n} P(s)$};

% ξ(s) ↔ RH、N(T) → RH：虚线标注（与 RH 相关的谱系连线）

% N(T) → RH：完成 ζ→N(T)→RH 线性逻辑链（虚线）
\draw[arr, eZero, cond, rounded corners=8pt]
  (NT.south west)
  to[out=-150, in=60, looseness=0.9]
  ($(rh.north)+(20pt, 0)$)
node[lab, pos=0.55, left=4pt, yshift=-430, xshift=-25, eZero, align=center]
    {$\mathrm{RH} \iff N_0(T)=N(T),\ \forall\,T>0$\\
   \tiny（$N_0(T)$：临界线 $\operatorname{Re}(s)=\dfrac{1}{2}$ 上虚部 $\leq T$ 的零点数）};

\draw[barr, eZero, cond]
  (xi.west) -- ($(rh.east)!0.0!(rh.south east)$)
  node[lab, midway, above=20pt,yshift=3pt,xshift=5, eZero, align=center]
    {$\mathrm{RH}\iff\xi\!\left(\dfrac{1}{2}+it\right)$\\
     的所有零点 $t$ 均为实数};

% M(x)→RH 箭头已删去：等价信息直接写在 RH 框内

% Gamma 自身解析延拓弧线：与表 tab:spectrum-stage8-guide 一致
\draw[arr, eExt, rounded corners=1.5pt]
  ($(gamma.north east)+(0,-5pt)$)
  to[out=0, in=0, looseness=1.6]
  ($(gamma.south east)+(0,5pt)$);
% 标签用独立节点锚定在 Γ 东侧弧线旁（仅递推式；延拓域已在框下半）
% 相对初版 (gamma.east+14pt)：净向左 3.7cm，上移 2pt
\node[lab, eExt, align=center, anchor=west]
  at ($(gamma.east)+({14pt-3.7cm},2pt)$)
  {递推延拓 $\Gamma(s)=\Gamma(s+1)/s$};

% Gamma（延拓后）→ ζ：Hankel 围道
\draw[arr, eExt]
  ($(gamma.south west)+(40pt,0)$)
  to[out=180, in=0, looseness=1.0]
  ($(zeta.south east)!0.45!(zeta.north east)$)
  node[lab, pos=0.5, yshift=-210, xshift=215, eExt, align=center]
    {$\displaystyle\zeta(s) = -\frac{\Gamma(1-s)}{2\pi i}\oint_{\mathrm{H}}\frac{(-z)^{s-1}}{e^z-1}\,\mathrm{d}z$\\
     解析延拓到 $\mathbb{C}\setminus\{1\}$，Hankel 围道消去分支切割};

\end{scope}% 主图

% ==================== 图例（取自 A-2 详细文案 + 分阶段灰白色） ====================
\node[draw=black!40, fill=gray!8, rounded corners=8pt, thick, align=center, inner sep=10pt, font=\tiny]
  (legend) at (0, -23.5) {
  \begin{tabular}{@{}ll@{\hspace{1.0cm}}ll@{}}
    \multicolumn{2}{l}{\scriptsize\textbf{节点颜色分类图例}} &
    \multicolumn{2}{l}{\scriptsize\textbf{知识谱系连线语义图例}} \\[6pt]
    \tikz[baseline=-0.5ex]\fill[lgPrime, draw=black!70, thick, rounded corners=1.5pt] (0,-0.12) rectangle (0.55,0.12); &
    \textcolor{lgTprime}{\textbf{素数计数函数}（实变，如 $\pi(x)$）} &
    \textcolor{lgEArith}{\rule[0.5ex]{16pt}{1.4pt}} &
    \textcolor{lgTarith}{\textbf{算术互逆}（M\"{o}bius / Stieltjes 反演）} \\[3pt]
    \tikz[baseline=-0.5ex]\fill[lgMu, draw=black!70, thick, rounded corners=1.5pt] (0,-0.12) rectangle (0.55,0.12); &
    \textcolor{lgTmu}{\textbf{算术积性函数}（离散，如 $\mu(n)$）} &
    \textcolor{lgEDef}{\rule[0.5ex]{16pt}{1.4pt}} &
    \textcolor{lgTdef}{\textbf{对象定义}（构造定义与主项提取）} \\[3pt]
    \tikz[baseline=-0.5ex]\fill[lgLi, draw=black!70, thick, rounded corners=1.5pt] (0,-0.12) rectangle (0.55,0.12); &
    \textcolor{lgTli}{\textbf{渐近逼近主项}（如 $\Li(x)$）} &
    \textcolor{lgEElem}{\rule[0.5ex]{16pt}{1.4pt}} &
    \textcolor{lgTelem}{\textbf{初等计数变换}（数论函数基本恒等式）} \\[3pt]
    \tikz[baseline=-0.5ex]\fill[lgXi, draw=black!55, thick, rounded corners=1.5pt] (0,-0.12) rectangle (0.55,0.12); &
    \textcolor{lgTxi}{\textbf{辅助复变函数}（如 $\xi(s)$）} &
    \textcolor{lgEZero}{\rule[0.5ex]{16pt}{1.4pt}} &
    \textcolor{lgTzero}{\textbf{零点结构与猜想}（分布特征及等价命题）} \\[3pt]
    \tikz[baseline=-0.5ex]\fill[lgRH, draw=black!55, thick, rounded corners=1.5pt] (0,-0.12) rectangle (0.55,0.12); &
    \textcolor{lgTrh}{\textbf{核心猜想与目标}（如 RH）} &
    \textcolor{lgEAnal}{\rule[0.5ex]{16pt}{1.4pt}} &
    \textcolor{lgTanal}{\textbf{分析桥梁}（积分变换与 Dirichlet 级数）} \\[3pt]
    \tikz[baseline=-0.5ex]{\clip[rounded corners=1.5pt] (0,-0.12) rectangle (0.55,0.12); \fill[lgGL] (0,-0.12) rectangle (0.55,0); \fill[lgGU] (0,0) rectangle (0.55,0.12); \draw[black!55, thick, rounded corners=1.5pt] (0,-0.12) rectangle (0.55,0.12); \draw[black!55, thick] (0,0) -- (0.55,0);} &
    \textcolor{lgTgamma}{\textbf{解析延拓复变}（初始域/延拓域）} &
    \textcolor{lgEPNT}{\rule[0.5ex]{16pt}{1.4pt}} &
    \textcolor{lgTpnt}{\textbf{渐近定理}（PNT 等确立的渐近推论）} \\[3pt]
    & & \textcolor{lgEExt}{\rule[0.5ex]{16pt}{1.4pt}} &
    \textcolor{lgText}{\textbf{解析延拓}（复围道积分与全纯延拓）} \\[3pt]
    & & \tikz[baseline=-0.5ex]{\draw[lgEZero, line width=1.2pt] (0.00,0) -- (0.12,0); \draw[lgEZero, line width=1.2pt] (0.17,0) -- (0.29,0); \draw[lgEZero, line width=1.2pt] (0.34,0) -- (0.46,0);} &
    \textbf{未证}
  \end{tabular}
};

\end{scope}% 位移：左 1.0cm、上 2.0cm

\end{tikzpicture}

% 几何与 黎曼猜想知识谱系关系图-A-2.tex 同步；未涉及者仅灰白化（含图例）
%
% \spectrumStage 仅为内部进度码（历史命名），不是章号：
%   2  → 第一次着色：第2章 初等算术层（fig:spectrum-stage2）
%   8  → 第二次着色：第9章 延拓/函数方程（fig:spectrum-stage8）
%  16  → 第三次着色：第15章 全文总图（fig:function-relations）
% 正文表述用「第一次 / 第二次 / 第三次着色」或章标签，勿写「第8章图」「第16章图」。
\providecommand{\spectrumStage}{16}
\makeatletter
\colorlet{cDim}{gray!16}\colorlet{dDim}{gray!48}\colorlet{tDim}{gray!52}
\ifnum\spectrumStage=2
  \colorlet{cPrime}{blue!12}\colorlet{dPrime}{black}\colorlet{tPrime}{black}
  \colorlet{cMu}{teal!15}\colorlet{dMu}{teal!60!black}\colorlet{tMu}{black}
  \colorlet{cLi}{yellow!25}\colorlet{dLi}{orange!70!black}\colorlet{tLi}{black}
  \colorlet{cGammaU}{gray!16}\colorlet{cGammaL}{gray!16}\colorlet{dGamma}{gray!48}\colorlet{tGamma}{gray!52}
  \colorlet{cZetaU}{gray!16}\colorlet{cZetaL}{gray!16}\colorlet{tZeta}{gray!52}
  \colorlet{cXi}{gray!16}\colorlet{dXi}{gray!48}\colorlet{tXi}{gray!52}
  \colorlet{cNT}{gray!16}\colorlet{dNT}{gray!48}\colorlet{tNT}{gray!52}
  \colorlet{cProd}{gray!16}\colorlet{dProd}{gray!48}\colorlet{tProd}{gray!52}
  \colorlet{cRH}{gray!16}\colorlet{dRH}{gray!48}\colorlet{tRH}{gray!52}
  \colorlet{eArith}{blue!70!black}\colorlet{eElem}{violet!80!black}\colorlet{eDef}{black!80}
  \colorlet{ePNT}{brown!70!black}\colorlet{eAnal}{gray!42}\colorlet{eExpl}{gray!42}
  \colorlet{eExt}{gray!42}\colorlet{eZero}{gray!42}\colorlet{eLate}{gray!42}
  \colorlet{lgPrime}{blue!12}\colorlet{lgMu}{teal!15}\colorlet{lgLi}{yellow!25}
  \colorlet{lgXi}{gray!22}\colorlet{lgRH}{gray!22}\colorlet{lgGU}{gray!22}\colorlet{lgGL}{gray!22}
  \colorlet{lgEArith}{blue!70!black}\colorlet{lgEDef}{black!80}\colorlet{lgEElem}{violet!80!black}
  \colorlet{lgEZero}{gray!45}\colorlet{lgEAnal}{gray!45}\colorlet{lgEPNT}{brown!70!black}\colorlet{lgEExt}{gray!45}
  \colorlet{lgTprime}{black}\colorlet{lgTmu}{black}\colorlet{lgTli}{black}
  \colorlet{lgTxi}{gray!50}\colorlet{lgTrh}{gray!50}\colorlet{lgTgamma}{gray!50}
  \colorlet{lgTarith}{black}\colorlet{lgTdef}{black}\colorlet{lgTelem}{black}
  \colorlet{lgTzero}{gray!50}\colorlet{lgTanal}{gray!50}\colorlet{lgTpnt}{black}\colorlet{lgText}{gray!50}
\else
\ifnum\spectrumStage=8
  \colorlet{cPrime}{blue!12}\colorlet{dPrime}{black}\colorlet{tPrime}{black}
  \colorlet{cMu}{teal!15}\colorlet{dMu}{teal!60!black}\colorlet{tMu}{black}
  \colorlet{cLi}{yellow!25}\colorlet{dLi}{orange!70!black}\colorlet{tLi}{black}
  \colorlet{cGammaU}{blue!15}\colorlet{cGammaL}{red!15}\colorlet{dGamma}{black}\colorlet{tGamma}{black}
  \colorlet{cZetaU}{blue!15}\colorlet{cZetaL}{red!15}\colorlet{tZeta}{black}
  \colorlet{cXi}{gray!16}\colorlet{dXi}{gray!48}\colorlet{tXi}{gray!52}
  \colorlet{cNT}{gray!16}\colorlet{dNT}{gray!48}\colorlet{tNT}{gray!52}
  \colorlet{cProd}{gray!16}\colorlet{dProd}{gray!48}\colorlet{tProd}{gray!52}
  \colorlet{cRH}{gray!16}\colorlet{dRH}{gray!48}\colorlet{tRH}{gray!52}
  \colorlet{eArith}{blue!70!black}\colorlet{eElem}{violet!80!black}\colorlet{eDef}{black!80}
  \colorlet{ePNT}{brown!70!black}\colorlet{eAnal}{green!50!black}\colorlet{eExpl}{gray!42}
  \colorlet{eExt}{red!70!black}\colorlet{eZero}{gray!42}\colorlet{eLate}{gray!42}
  \colorlet{lgPrime}{blue!12}\colorlet{lgMu}{teal!15}\colorlet{lgLi}{yellow!25}
  \colorlet{lgXi}{gray!22}\colorlet{lgRH}{gray!22}\colorlet{lgGU}{blue!15}\colorlet{lgGL}{red!15}
  \colorlet{lgEArith}{blue!70!black}\colorlet{lgEDef}{black!80}\colorlet{lgEElem}{violet!80!black}
  \colorlet{lgEZero}{gray!45}\colorlet{lgEAnal}{green!50!black}\colorlet{lgEPNT}{brown!70!black}\colorlet{lgEExt}{red!70!black}
  \colorlet{lgTprime}{black}\colorlet{lgTmu}{black}\colorlet{lgTli}{black}
  \colorlet{lgTxi}{gray!50}\colorlet{lgTrh}{gray!50}\colorlet{lgTgamma}{black}
  \colorlet{lgTarith}{black}\colorlet{lgTdef}{black}\colorlet{lgTelem}{black}
  \colorlet{lgTzero}{gray!50}\colorlet{lgTanal}{black}\colorlet{lgTpnt}{black}\colorlet{lgText}{black}
\else
  \colorlet{cPrime}{blue!12}\colorlet{dPrime}{black}\colorlet{tPrime}{black}
  \colorlet{cMu}{teal!15}\colorlet{dMu}{teal!60!black}\colorlet{tMu}{black}
  \colorlet{cLi}{yellow!25}\colorlet{dLi}{orange!70!black}\colorlet{tLi}{black}
  \colorlet{cGammaU}{blue!15}\colorlet{cGammaL}{red!15}\colorlet{dGamma}{black}\colorlet{tGamma}{black}
  \colorlet{cZetaU}{blue!15}\colorlet{cZetaL}{red!15}\colorlet{tZeta}{black}
  \colorlet{cXi}{cyan!15}\colorlet{dXi}{cyan!60!black}\colorlet{tXi}{black}
  \colorlet{cNT}{cyan!15}\colorlet{dNT}{cyan!60!black}\colorlet{tNT}{black}
  \colorlet{cProd}{cyan!15}\colorlet{dProd}{cyan!60!black}\colorlet{tProd}{black}
  \colorlet{cRH}{orange!15}\colorlet{dRH}{orange!70!black}\colorlet{tRH}{black}
  \colorlet{eArith}{blue!70!black}\colorlet{eElem}{violet!80!black}\colorlet{eDef}{black!80}
  \colorlet{ePNT}{brown!70!black}\colorlet{eAnal}{green!50!black}\colorlet{eExpl}{green!50!black}
  \colorlet{eExt}{red!70!black}\colorlet{eZero}{orange!80!black}\colorlet{eLate}{black!80}
  \colorlet{lgPrime}{blue!12}\colorlet{lgMu}{teal!15}\colorlet{lgLi}{yellow!25}
  \colorlet{lgXi}{cyan!15}\colorlet{lgRH}{orange!15}\colorlet{lgGU}{blue!15}\colorlet{lgGL}{red!15}
  \colorlet{lgEArith}{blue!70!black}\colorlet{lgEDef}{black!80}\colorlet{lgEElem}{violet!80!black}
  \colorlet{lgEZero}{orange!80!black}\colorlet{lgEAnal}{green!50!black}\colorlet{lgEPNT}{brown!70!black}\colorlet{lgEExt}{red!70!black}
  \colorlet{lgTprime}{black}\colorlet{lgTmu}{black}\colorlet{lgTli}{black}
  \colorlet{lgTxi}{black}\colorlet{lgTrh}{black}\colorlet{lgTgamma}{black}
  \colorlet{lgTarith}{black}\colorlet{lgTdef}{black}\colorlet{lgTelem}{black}
  \colorlet{lgTzero}{black}\colorlet{lgTanal}{black}\colorlet{lgTpnt}{black}\colorlet{lgText}{black}
\fi\fi
\makeatother

\begin{tikzpicture}[scale=1.0, transform shape,
  box/.style    = {draw, rounded corners=5pt, minimum width=2.8cm,
                   minimum height=0.9cm, align=center, font=\small, thick},
  zbox/.style   = {draw, rounded corners=5pt, minimum width=2.8cm,
                   minimum height=0.9cm, align=center, font=\small,
                   thick, fill=gray!18},
  splitbox/.style = {draw, rounded corners=5pt, minimum width=7cm,
                     minimum height=2.2cm, align=center, font=\small, thick,
                     rectangle split, rectangle split parts=2,
                     rectangle split part fill={red!8, white}},
  arr/.style    = {-{Stealth[length=6pt,width=5pt]}, thick},
  barr/.style   = {{Stealth[length=6pt,width=5pt]}-{Stealth[length=6pt,width=5pt]}, thick},
  % 线型：虚线仅用于 N(T)→RH、ξ↔RH；其余为实线
  proven/.style = {solid},
  cond/.style   = {dashed},
  lab/.style    = {font=\tiny, align=center, inner sep=3pt},
  rhbox/.style  = {draw=dRH, rounded corners=5pt, minimum width=8cm,
                   minimum height=1.0cm, align=center, font=\small,
                   fill=cRH, text=tRH, thick}
]

% 布局与 A-2 同步：左移 1.0cm、上移 2.0cm
% 左右略放宽，避免 Γ/ζ 自环旁公式标注被裁切
\path[use as bounding box]
  (-10.2, -26.2) rectangle (10.0, 7.4);

\begin{scope}[shift={(-1.0,2.0)}]

\node[font=\bfseries\Large, align=center] at (0, 4.35)
  {黎曼猜想知识谱系关系图};

% 主图单独上移 1cm；标题与图例位置不动
\begin{scope}[shift={(-2.75,-1.2)}]
%% ==================== 节点分层 ====================

% 第一层：素数计数函数 (y=0)
\node[box, fill=cPrime, draw=dPrime]    (J)    at (-4.2,  0.0) {$J(x)$};
\node[box, fill=cPrime, draw=dPrime]    (pi)   at ( 0.0,  0.0) {$\pi(x)$};
\node[box, fill=cPrime, draw=dPrime]    (th)   at ( 5.0,  0.0) {$\theta(x)$};
\node[box, fill=cPrime, draw=dPrime]    (ps)   at ( 9.7,  0.0) {$\psi(x)$};

% 第二层：积性函数 (y=-3.0)
% 节点定义修改
\node[box, fill=cMu, draw=dMu] (mu)  at (0.3, -3.0) {$\mu(n)$};
\node[box, fill=cMu, draw=dMu] (Lam) at (5.0, -3.0) {$\Lambda(n)$};

%Weierstrass乘积
\node[box, fill=cProd, draw=dProd, text=tProd]  (prod)   at ( 9.5,  -17.0) {\tiny\text{Hadamard 乘积}};
% M(x) 节点已移除：初等等价信息直接标注在 RH 框内

% 第三层：Li 与 Gamma (y=-6.0 至 -7.0)
\node[box, fill=cLi, draw=dLi] (Li) at (-2.0, -6.0) {$\Li(x)$};

% Gamma(s) 上下分色：上部标收敛域，下部标延拓
\node[splitbox, rectangle split draw splits=false, minimum width=5cm,
      minimum height=4.7cm,
      rectangle split part fill={cGammaU, cGammaL}] (gamma) at (7.0, -6.0)
    {
     \centering \color{tGamma}$\operatorname{Re}(s)>0$
     \nodepart{two}
     \raisebox{6pt}{\color{tGamma}$\mathbb{C}\setminus\{0,-1,-2,\dots\}$}
    };
% Gamma(s) 标注（固定在节点中线左侧）
\node[font=\small, left=-30pt, text=tGamma] at (gamma.west) {$\Gamma(s)$};

% 【修订3】ζ(s) 节点：下半部分补全函数方程左侧 ζ(s)=
\node[splitbox, rectangle split draw splits=false, minimum height=4cm,
      rectangle split part fill={cZetaU, cZetaL}] (zeta) at (0.3, -10.5)
    {\centering\small
     \color{tZeta}$\displaystyle
       \sum_{n=1}^{\infty}\frac{1}{n^s}
       =\prod_{p}\dfrac{1}{1-p^{-s}}
       \quad(\operatorname{Re}(s)>1)$%
     \vphantom{$\displaystyle\sum_{n=1}^{\infty}$}%
     \nodepart{two}
     % 修订：补全 ζ(s)= 使函数方程完整
     \textcolor{tZeta}{$\displaystyle\zeta(s)=
       2^s\pi^{s-1}
       \sin\!\left(\frac{\pi s}{2}\right)\Gamma(1-s)\,\zeta(1-s)$%
     \vphantom{$\displaystyle\sum_{n=1}^{\infty}\mathbb{C}\setminus\{1\}$}}
    };
% ζ(s) 标注


% ξ(s) 节点：下移至与 RH 同一水平，右侧对称于 M(x)
\node[box, fill=cXi, draw=dXi, text=tXi, minimum height=1.2cm]
  (xi) at (4.85, -17.0)
  {$\xi(s)$\\ \tiny 非平凡零点$\rho$与$\zeta(s)$相同\\
  \tiny $\xi(s)=\xi(1-s)$(函数方程推论)
  };


% 【修订8】N(T) 节点：置于 ζ 与 RH 正中间，构成竖直逻辑链
\node[box, fill=cNT, draw=dNT, text=tNT, minimum height=1.3cm]
  (NT) at (2.3, -14.0)
  {$N(T)$\\ \tiny $\zeta(s)$ 零点计数函数\\[1pt]\tiny（幅角原理）};

% ==================== 黎曼猜想节点 ====================
% RH 下移至 y=-20.5，为 N(T) 留出充分空间
\node[rhbox, minimum width=4.6cm, minimum height=2.2cm] (rh) at (-1.9, -17.0) {
    {\tiny\color{tRH} 黎曼猜想 (Riemann Hypothesis, 1859)}\\[3pt]
    {\tiny\color{tRH}$\zeta(s)$ 的所有非平凡零点 $\rho$ 满足
    $\operatorname{Re}(\rho)=\dfrac{1}{2}$}\\[4pt]
{\tiny\color{tRH} 初等等价（梅滕斯函数）：%
  $\displaystyle M(x)=\sum_{n\leq x}\mu(n)$}\\
{\tiny\color{tRH}
  $\mathrm{RH} \iff M(x) = O(x^{1/2+\varepsilon}),\ \forall \varepsilon>0$}
};


% ==================== 连线 ====================

% ----- 第一层内部 -----
\draw[barr, eArith] (J.east) -- (pi.west)
  node[lab, midway, above=14pt, xshift=-2pt, eArith]
    {$\displaystyle\pi(x)=\sum_{n=1}^{\infty}\frac{\mu(n)}{n}J(x^{1/n})$}
  node[lab, midway, below=14pt, xshift=-2pt, eArith]
    {$\displaystyle J(x)=\sum_{n=1}^{\infty}\frac{1}{n}\pi(x^{1/n})$};

\draw[barr, eElem] (pi.east) -- (th.west)
  node[lab, midway, above=14pt, eElem]
    {$\displaystyle\theta(x)=\pi(x)\log x-\int_2^x\frac{\pi(t)}{t}\,\mathrm{d}t$}
  node[lab, midway, above={14pt}, yshift={-50pt}, eElem]
    {\tiny $\displaystyle\pi(x) = \dfrac{\theta(x)}{\log x}+ \int_{2}^{x} \frac{\theta(t)}{t\,(\log t)^2} \,\mathrm{d}t$};

\draw[barr, eElem] (th.east) -- (ps.west)
  node[lab, midway, above=11pt, xshift={5pt}, eElem]
    {$\displaystyle\psi(x)=\sum_{k=1}^{\infty}\theta(x^{1/k})$}
  node[lab, midway, above=-31pt, xshift={5pt}, yshift={-8pt}, eElem]
    {\tiny $\displaystyle\theta(x)=\sum_{k=1}^{\infty}\mu(k)\psi(x^{1/k})$};


% ----- 第二层内部 -----
\draw[barr, eArith] (mu.east) -- (Lam.west)
  node[lab, midway, yshift=20, xshift=0.3cm, eArith]
    {$\displaystyle\Lambda(n)=-\sum_{d\mid n}\mu(d)\log d$
     \quad{\rm (M\"{o}bius 反演)}}
  node[lab, midway, yshift={-12pt-0.5cm}, xshift=0, eArith]
    {$\displaystyle\log n=\sum_{d\mid n}\Lambda(d)$};

% μ(n) → M(x) 的定义连线已省去（见节点内文字），避免与已有连线交叉

% Lambda -> psi
\draw[arr, eDef] (Lam.east) to[out=0, in=-90, looseness=1.0]
  ($(ps.south west)+(30pt,0)$)
  node[lab, pos=0.5, yshift=-90, xshift=250, eDef]
    {$\displaystyle\psi(x)=\sum_{n\le x}\Lambda(n)$};

% ----- 第三层：Li 和 Gamma -----

% 【修订5】Li(x) → π(x)：改为渐近等价语义（非因果推导）
\draw[arr, ePNT]
  ($(Li.north)!.75!(Li.north east)$)
  to[out=90, in=-135, looseness=1.8]
  ($(pi.south west)+(35pt,0)$)
  node[lab, pos=0.4, yshift=-45, xshift=45, ePNT, align=left]
    {\textbf{素数定理（PNT）}：$\pi(x)\sim\Li(x)$\\
     (若RH成立，$\pi(x)=\Li(x)+O(\sqrt{x}\log x)$)};

% Li -> J（显式公式主项；属 eExpl：第2/9章灰白，第15章全文总图着色）
\draw[arr, eExpl]
  (Li.north) to[out=135, in=-90, looseness=1.5]
  ($(J.south east)+(-22pt,0)$)
  node[lab, pos=0.5, yshift=-115, xshift=-65, eExpl, align=center]
    {$\Li(x)$ 是 $J(x)$ 显式\\[2pt]公式的主项};

% Gamma（Re>0）→ ζ：Euler 积分表示（该积分公式本身要求 Re(s)>1）
\draw[arr, eAnal]
  ([yshift=6pt] gamma.west)
  to[out=180, in=90, looseness=0.8]
  ($(zeta.north)+(30pt,0)$)
  node[lab, pos=0.5, yshift=-235, xshift=90,eAnal, align=center]
    {$\displaystyle\zeta(s)\Gamma(s)
      =\int_0^{\infty}\frac{x^{s-1}}{e^x-1}\,\mathrm{d}x$\\
     $(\operatorname{Re}(s)>1)$};

% μ(n) → ζ(s)
\draw[arr, eAnal]
  (mu.south) to[out=-90, in=90, looseness=1.2]
  ($(zeta.north)+({-20pt+0.3cm},0)$)
  node[lab, pos=0.5, yshift={-120pt-0.5cm-0.4cm}, xshift={45pt-5pt}, eAnal]
    {$\displaystyle\frac{1}{\zeta(s)}=\sum_{n=1}^{\infty}\frac{\mu(n)}{n^s}$\\
     $(\operatorname{Re}(s)>1)$};

% ----- J(x) 与 ψ(x) 的 Stieltjes 积分互逆关系 -----
% 【修订4】加注 dψ(t)/ln t 在 t=1 附近需 t>1 的说明
\draw[barr, eArith, line width=0.8pt, rounded corners=8pt]
  ($(J.north)+(0,0.0cm)$)
  -- ++(0,2.0cm)
  -- ++(13.5cm,0)
  -- ++(0,-2.0cm)
  node[lab, pos=0.35, above=6pt, yshift={10pt+0.5cm}, xshift=-200, eArith, align=center]
    {$\displaystyle\psi(x)=\int_{2^-}^x\log t\,\mathrm{d}J(t)$\quad(Stieltjes 积分，$J\to\psi$)}
  node[lab, pos=0.65, below=6pt, yshift={32pt+0.2cm}, xshift=-205, eArith, align=center]
    {$\displaystyle J(x)=\int_{2^-}^x\frac{\mathrm{d}\psi(t)}{\log t}$
     \quad(Stieltjes 反演，$\psi\to J$；积分自 $t>1$)};

% ----- 第四层：ζ 和 ξ 与第二层的连线 -----

% ζ → ψ：显式公式（Perron 积分 + 留数，单向；x≠p^k 时公式精确成立）
\draw[arr, eExpl, rounded corners=30pt]
  ($(zeta.east)+(0, -10pt)$) -| ($(ps.south)+(30pt,0)$)
  node[lab, pos=0.85, right=10pt, yshift=-170, xshift={-170pt-0.2cm},eExpl, align=left]
{$\psi(x)$ 显式公式（Perron 积分，$\zeta\to\psi$）\\
     $\displaystyle \psi(x) = \frac{1}{2\pi i}\int_{c-i\infty}^{c+i\infty}\!\left(-\frac{\zeta'(s)}{\zeta(s)}\right)\!\frac{x^s}{s}\,\mathrm{d}s$\\
     $\displaystyle = x-\sum_{\rho}\frac{x^{\rho}}{\rho}-\log(2\pi)-\tfrac{1}{2}\log(1-x^{-2})$\\
     \tiny $x>1,\;x\neq p^k$；$\rho$ 按 $|\mathrm{Im}(\rho)|$ 升序配对，条件收敛};

% ψ → ζ：Mellin 变换（单向）
\draw[arr, eExpl, rounded corners=30pt]
  ($(ps.south)+(40pt,0)$) |- ($(zeta.east)+(0, -20pt)$)
  node[lab, pos=0.8, yshift=-27, xshift=30,eExpl, align=center]
    {Mellin 变换（$\psi(x) \longleftrightarrow -\dfrac{\zeta'(s)}{\zeta(s)}$，即负对数导数）\\
     $\displaystyle -\frac{\zeta'(s)}{\zeta(s)} = s\int_{1}^{\infty}\frac{\psi(x)}{x^{s+1}}\,\mathrm{d}x$};
       
       
       
       
       
       
       
% Lambda -> zeta（对数微分 Dirichlet 级数）
\draw[arr, eAnal]
  ($(Lam.south west)+(6pt,0)$) to[out=-90, in=90, looseness=0.4]
  ($(zeta.north west)!0.50!(zeta.north east)$)
  node[lab, pos=0.5, yshift={-120pt-0.5cm-0.4cm+1cm}, xshift={145pt-5pt}, eAnal]
    {$\displaystyle-\frac{\zeta'(s)}{\zeta(s)}
      =\sum_{n=1}^{\infty}\frac{\Lambda(n)}{n^s}$};

% J(x) 显式公式（ζ → J）
\draw[arr, eExpl, rounded corners=50pt]
  ($(zeta.south west)+(0,6pt)$) -- ++(-60pt,0) -- ($(J.south west)+(5pt,0)$)
  node[lab, pos=0.5, yshift=-203, xshift={75pt-0.2cm}, eExpl, align=center]
{\textbf{Perron 公式} $\Rightarrow$ $J(x)$ 显式展开\\
     $\displaystyle
     \begin{aligned}
         &J(x)=\frac{1}{2\pi i}\int_{c-i\infty}^{c+i\infty}
           \frac{\log\zeta(s)}{s}\,x^s\,\mathrm{d}s\\
         &=\Li(x)-\sum_{\rho}\Li(x^{\rho})
           -\log 2+\int_x^{\infty}\frac{\mathrm{d}t}{t(t^2-1)\log t}
     \end{aligned}$\\
     \tiny $x>1,\;x\neq p^k$；$\rho$ 按 $|\mathrm{Im}(\rho)|$ 升序配对，条件收敛};

% ζ → J：Mellin 变换（ln ζ）
\draw[arr, eAnal, rounded corners=5pt]
  (zeta.north west) to[out=110, in=-90, looseness=1.5] (J.south)
  node[lab, pos=0.4, yshift=-205, xshift=-55, eAnal, align=center]
    {$\displaystyle\log\zeta(s)=s\int_1^{\infty}J(x)\,x^{-s-1}\,\mathrm{d}x$\\
     $(\operatorname{Re}(s)>1)$};

% ζ → ξ（定义）：ξ 已下移，连线沿右侧空白区域向右下延伸
\draw[arr, eLate, rounded corners=15pt]
  ($(zeta.south east)!0.15!(zeta.north east)$)
  -- ++(1.0, 0)
  -- ++(0, -3.8)
  -- ($(xi.north)+(0pt,0)$)
  node[lab, pos=0.25, right=4pt,yshift=30, xshift=0, eLate, align=center]
    {$\xi(s)=\dfrac{1}{2}s(s-1)\pi^{-s/2}\Gamma(s/2)\zeta(s)$  
};






% xi -> prod (Weierstrass product connection)
\draw[arr, eLate]
  %(xi.south) -- (prod.south)
  (xi.east) --(prod.west)
  node[midway, right=4pt, yshift=28, xshift=-45, eLate]
    {\tiny$\displaystyle\xi(s)=e^{A+Bs}\prod_{\rho}\left(1-\frac{s}{\rho}\right)e^{s/\rho}$};

% prod -> ps (zero sum in explicit formula)
%\draw[arr, violet!70!black, rounded corners=15pt]
  %(prod.east) -- ++(4.4, 0) |- (ps.south)
  %node[lab, pos=0.6, right=2pt, violet!70!black, align=left]
   % {对数微分与 Perron 积分\\
    % 给出显式公式中\\
    % 零点和 $\sum_{\rho}\frac{x_{\rho}}{\rho}$ 项};
%
% ζ 自身左侧弧线（解析延拓）：与表 tab:spectrum-stage8-guide 一致
\draw[arr, eExt, rounded corners=5pt]
  ($(zeta.north west)!0.15!(zeta.south west)$)
  to[out=180, in=180, looseness=1.6]
  ($(zeta.north west)!0.85!(zeta.south west)$);
% 标签用独立节点锚定在 ζ 西侧弧线旁（避免 path 末 node 漂移）
% 单行；中心在 zeta.west 右 2.5cm（相对初版再累计右移）
\node[lab, eExt, align=center, anchor=center]
  at ([xshift=2.5cm]zeta.west)
  {$\zeta$ 亚纯延拓至 $\mathbb{C}\setminus\{1\}$（Hankel / $\vartheta$ 两法）};

% ζ → N(T)：竖直向下（弧线）
\draw[arr, eZero]
  ($(zeta.south east)!0.10!(zeta.north east)$)
  to[out=0, in=90] ++(0.8, -0.8)
  to[out=-90, in=0] (NT.east)
  node[lab, pos=0.6, right=6pt, yshift=-357, xshift=5, eZero, align=center]
    {$\displaystyle N(T)=\frac{T}{2\pi}\log\frac{T}{2\pi e}+O(\log T)$\\
     \tiny 零点计数公式（von Mangoldt）};

% 【已删去】ζ → RH 蓝色直连线（设计意图改为ζ→N(T)→RH线性链）
%\draw[arr, eArith, rounded corners=10pt]
%  (zeta.south)
%  -- ++(0, -3.35)
%  -- ++(-4.1, 0)
%  -- ($(rh.north)!0.80!(rh.north west)$)
%  node[lab, pos=0.35, left=6pt, yshift=20pt,xshift=120,eArith, align=center]
%    {
%     $\zeta(s)$ 的非平凡零点全部位于直线 $\operatorname{Re}(s)=\dfrac12$ 上};
% RH -- P(s) 画线 + 文字
%\draw[arr, cyan!60!black]
 % ($(rh.east)!0.60!(rh.south east)$) -- (prod.west)    % 起点 -- 终点（必须有！）
  %node[midway, right=6pt, font=\tiny, cyan!60!black, align=center]
  %  {$\zeta(s) = \frac{e^{s(\log(2\pi) - 1)}}{2(s-1)} \prod_{n=1}^{\infty} \left(1 + \frac{s}{2n}\right)e^{-s/2n} P(s)$};

% ξ(s) ↔ RH、N(T) → RH：虚线标注（与 RH 相关的谱系连线）

% N(T) → RH：完成 ζ→N(T)→RH 线性逻辑链（虚线）
\draw[arr, eZero, cond, rounded corners=8pt]
  (NT.south west)
  to[out=-150, in=60, looseness=0.9]
  ($(rh.north)+(20pt, 0)$)
node[lab, pos=0.55, left=4pt, yshift=-430, xshift=-25, eZero, align=center]
    {$\mathrm{RH} \iff N_0(T)=N(T),\ \forall\,T>0$\\
   \tiny（$N_0(T)$：临界线 $\operatorname{Re}(s)=\dfrac{1}{2}$ 上虚部 $\leq T$ 的零点数）};

\draw[barr, eZero, cond]
  (xi.west) -- ($(rh.east)!0.0!(rh.south east)$)
  node[lab, midway, above=20pt,yshift=3pt,xshift=5, eZero, align=center]
    {$\mathrm{RH}\iff\xi\!\left(\dfrac{1}{2}+it\right)$\\
     的所有零点 $t$ 均为实数};

% M(x)→RH 箭头已删去：等价信息直接写在 RH 框内

% Gamma 自身解析延拓弧线：与表 tab:spectrum-stage8-guide 一致
\draw[arr, eExt, rounded corners=1.5pt]
  ($(gamma.north east)+(0,-5pt)$)
  to[out=0, in=0, looseness=1.6]
  ($(gamma.south east)+(0,5pt)$);
% 标签用独立节点锚定在 Γ 东侧弧线旁（仅递推式；延拓域已在框下半）
% 相对初版 (gamma.east+14pt)：净向左 3.7cm，上移 2pt
\node[lab, eExt, align=center, anchor=west]
  at ($(gamma.east)+({14pt-3.7cm},2pt)$)
  {递推延拓 $\Gamma(s)=\Gamma(s+1)/s$};

% Gamma（延拓后）→ ζ：Hankel 围道
\draw[arr, eExt]
  ($(gamma.south west)+(40pt,0)$)
  to[out=180, in=0, looseness=1.0]
  ($(zeta.south east)!0.45!(zeta.north east)$)
  node[lab, pos=0.5, yshift=-210, xshift=215, eExt, align=center]
    {$\displaystyle\zeta(s) = -\frac{\Gamma(1-s)}{2\pi i}\oint_{\mathrm{H}}\frac{(-z)^{s-1}}{e^z-1}\,\mathrm{d}z$\\
     解析延拓到 $\mathbb{C}\setminus\{1\}$，Hankel 围道消去分支切割};

\end{scope}% 主图

% ==================== 图例（取自 A-2 详细文案 + 分阶段灰白色） ====================
\node[draw=black!40, fill=gray!8, rounded corners=8pt, thick, align=center, inner sep=10pt, font=\tiny]
  (legend) at (0, -23.5) {
  \begin{tabular}{@{}ll@{\hspace{1.0cm}}ll@{}}
    \multicolumn{2}{l}{\scriptsize\textbf{节点颜色分类图例}} &
    \multicolumn{2}{l}{\scriptsize\textbf{知识谱系连线语义图例}} \\[6pt]
    \tikz[baseline=-0.5ex]\fill[lgPrime, draw=black!70, thick, rounded corners=1.5pt] (0,-0.12) rectangle (0.55,0.12); &
    \textcolor{lgTprime}{\textbf{素数计数函数}（实变，如 $\pi(x)$）} &
    \textcolor{lgEArith}{\rule[0.5ex]{16pt}{1.4pt}} &
    \textcolor{lgTarith}{\textbf{算术互逆}（M\"{o}bius / Stieltjes 反演）} \\[3pt]
    \tikz[baseline=-0.5ex]\fill[lgMu, draw=black!70, thick, rounded corners=1.5pt] (0,-0.12) rectangle (0.55,0.12); &
    \textcolor{lgTmu}{\textbf{算术积性函数}（离散，如 $\mu(n)$）} &
    \textcolor{lgEDef}{\rule[0.5ex]{16pt}{1.4pt}} &
    \textcolor{lgTdef}{\textbf{对象定义}（构造定义与主项提取）} \\[3pt]
    \tikz[baseline=-0.5ex]\fill[lgLi, draw=black!70, thick, rounded corners=1.5pt] (0,-0.12) rectangle (0.55,0.12); &
    \textcolor{lgTli}{\textbf{渐近逼近主项}（如 $\Li(x)$）} &
    \textcolor{lgEElem}{\rule[0.5ex]{16pt}{1.4pt}} &
    \textcolor{lgTelem}{\textbf{初等计数变换}（数论函数基本恒等式）} \\[3pt]
    \tikz[baseline=-0.5ex]\fill[lgXi, draw=black!55, thick, rounded corners=1.5pt] (0,-0.12) rectangle (0.55,0.12); &
    \textcolor{lgTxi}{\textbf{辅助复变函数}（如 $\xi(s)$）} &
    \textcolor{lgEZero}{\rule[0.5ex]{16pt}{1.4pt}} &
    \textcolor{lgTzero}{\textbf{零点结构与猜想}（分布特征及等价命题）} \\[3pt]
    \tikz[baseline=-0.5ex]\fill[lgRH, draw=black!55, thick, rounded corners=1.5pt] (0,-0.12) rectangle (0.55,0.12); &
    \textcolor{lgTrh}{\textbf{核心猜想与目标}（如 RH）} &
    \textcolor{lgEAnal}{\rule[0.5ex]{16pt}{1.4pt}} &
    \textcolor{lgTanal}{\textbf{分析桥梁}（积分变换与 Dirichlet 级数）} \\[3pt]
    \tikz[baseline=-0.5ex]{\clip[rounded corners=1.5pt] (0,-0.12) rectangle (0.55,0.12); \fill[lgGL] (0,-0.12) rectangle (0.55,0); \fill[lgGU] (0,0) rectangle (0.55,0.12); \draw[black!55, thick, rounded corners=1.5pt] (0,-0.12) rectangle (0.55,0.12); \draw[black!55, thick] (0,0) -- (0.55,0);} &
    \textcolor{lgTgamma}{\textbf{解析延拓复变}（初始域/延拓域）} &
    \textcolor{lgEPNT}{\rule[0.5ex]{16pt}{1.4pt}} &
    \textcolor{lgTpnt}{\textbf{渐近定理}（PNT 等确立的渐近推论）} \\[3pt]
    & & \textcolor{lgEExt}{\rule[0.5ex]{16pt}{1.4pt}} &
    \textcolor{lgText}{\textbf{解析延拓}（复围道积分与全纯延拓）} \\[3pt]
    & & \tikz[baseline=-0.5ex]{\draw[lgEZero, line width=1.2pt] (0.00,0) -- (0.12,0); \draw[lgEZero, line width=1.2pt] (0.17,0) -- (0.29,0); \draw[lgEZero, line width=1.2pt] (0.34,0) -- (0.46,0);} &
    \textbf{未证}
  \end{tabular}
};

\end{scope}% 位移：左 1.0cm、上 2.0cm

\end{tikzpicture}

% 几何与 黎曼猜想知识谱系关系图-A-2.tex 同步；未涉及者仅灰白化（含图例）
%
% \spectrumStage 仅为内部进度码（历史命名），不是章号：
%   2  → 第一次着色：第2章 初等算术层（fig:spectrum-stage2）
%   8  → 第二次着色：第9章 延拓/函数方程（fig:spectrum-stage8）
%  16  → 第三次着色：第15章 全文总图（fig:function-relations）
% 正文表述用「第一次 / 第二次 / 第三次着色」或章标签，勿写「第8章图」「第16章图」。
\providecommand{\spectrumStage}{16}
\makeatletter
\colorlet{cDim}{gray!16}\colorlet{dDim}{gray!48}\colorlet{tDim}{gray!52}
\ifnum\spectrumStage=2
  \colorlet{cPrime}{blue!12}\colorlet{dPrime}{black}\colorlet{tPrime}{black}
  \colorlet{cMu}{teal!15}\colorlet{dMu}{teal!60!black}\colorlet{tMu}{black}
  \colorlet{cLi}{yellow!25}\colorlet{dLi}{orange!70!black}\colorlet{tLi}{black}
  \colorlet{cGammaU}{gray!16}\colorlet{cGammaL}{gray!16}\colorlet{dGamma}{gray!48}\colorlet{tGamma}{gray!52}
  \colorlet{cZetaU}{gray!16}\colorlet{cZetaL}{gray!16}\colorlet{tZeta}{gray!52}
  \colorlet{cXi}{gray!16}\colorlet{dXi}{gray!48}\colorlet{tXi}{gray!52}
  \colorlet{cNT}{gray!16}\colorlet{dNT}{gray!48}\colorlet{tNT}{gray!52}
  \colorlet{cProd}{gray!16}\colorlet{dProd}{gray!48}\colorlet{tProd}{gray!52}
  \colorlet{cRH}{gray!16}\colorlet{dRH}{gray!48}\colorlet{tRH}{gray!52}
  \colorlet{eArith}{blue!70!black}\colorlet{eElem}{violet!80!black}\colorlet{eDef}{black!80}
  \colorlet{ePNT}{brown!70!black}\colorlet{eAnal}{gray!42}\colorlet{eExpl}{gray!42}
  \colorlet{eExt}{gray!42}\colorlet{eZero}{gray!42}\colorlet{eLate}{gray!42}
  \colorlet{lgPrime}{blue!12}\colorlet{lgMu}{teal!15}\colorlet{lgLi}{yellow!25}
  \colorlet{lgXi}{gray!22}\colorlet{lgRH}{gray!22}\colorlet{lgGU}{gray!22}\colorlet{lgGL}{gray!22}
  \colorlet{lgEArith}{blue!70!black}\colorlet{lgEDef}{black!80}\colorlet{lgEElem}{violet!80!black}
  \colorlet{lgEZero}{gray!45}\colorlet{lgEAnal}{gray!45}\colorlet{lgEPNT}{brown!70!black}\colorlet{lgEExt}{gray!45}
  \colorlet{lgTprime}{black}\colorlet{lgTmu}{black}\colorlet{lgTli}{black}
  \colorlet{lgTxi}{gray!50}\colorlet{lgTrh}{gray!50}\colorlet{lgTgamma}{gray!50}
  \colorlet{lgTarith}{black}\colorlet{lgTdef}{black}\colorlet{lgTelem}{black}
  \colorlet{lgTzero}{gray!50}\colorlet{lgTanal}{gray!50}\colorlet{lgTpnt}{black}\colorlet{lgText}{gray!50}
\else
\ifnum\spectrumStage=8
  \colorlet{cPrime}{blue!12}\colorlet{dPrime}{black}\colorlet{tPrime}{black}
  \colorlet{cMu}{teal!15}\colorlet{dMu}{teal!60!black}\colorlet{tMu}{black}
  \colorlet{cLi}{yellow!25}\colorlet{dLi}{orange!70!black}\colorlet{tLi}{black}
  \colorlet{cGammaU}{blue!15}\colorlet{cGammaL}{red!15}\colorlet{dGamma}{black}\colorlet{tGamma}{black}
  \colorlet{cZetaU}{blue!15}\colorlet{cZetaL}{red!15}\colorlet{tZeta}{black}
  \colorlet{cXi}{gray!16}\colorlet{dXi}{gray!48}\colorlet{tXi}{gray!52}
  \colorlet{cNT}{gray!16}\colorlet{dNT}{gray!48}\colorlet{tNT}{gray!52}
  \colorlet{cProd}{gray!16}\colorlet{dProd}{gray!48}\colorlet{tProd}{gray!52}
  \colorlet{cRH}{gray!16}\colorlet{dRH}{gray!48}\colorlet{tRH}{gray!52}
  \colorlet{eArith}{blue!70!black}\colorlet{eElem}{violet!80!black}\colorlet{eDef}{black!80}
  \colorlet{ePNT}{brown!70!black}\colorlet{eAnal}{green!50!black}\colorlet{eExpl}{gray!42}
  \colorlet{eExt}{red!70!black}\colorlet{eZero}{gray!42}\colorlet{eLate}{gray!42}
  \colorlet{lgPrime}{blue!12}\colorlet{lgMu}{teal!15}\colorlet{lgLi}{yellow!25}
  \colorlet{lgXi}{gray!22}\colorlet{lgRH}{gray!22}\colorlet{lgGU}{blue!15}\colorlet{lgGL}{red!15}
  \colorlet{lgEArith}{blue!70!black}\colorlet{lgEDef}{black!80}\colorlet{lgEElem}{violet!80!black}
  \colorlet{lgEZero}{gray!45}\colorlet{lgEAnal}{green!50!black}\colorlet{lgEPNT}{brown!70!black}\colorlet{lgEExt}{red!70!black}
  \colorlet{lgTprime}{black}\colorlet{lgTmu}{black}\colorlet{lgTli}{black}
  \colorlet{lgTxi}{gray!50}\colorlet{lgTrh}{gray!50}\colorlet{lgTgamma}{black}
  \colorlet{lgTarith}{black}\colorlet{lgTdef}{black}\colorlet{lgTelem}{black}
  \colorlet{lgTzero}{gray!50}\colorlet{lgTanal}{black}\colorlet{lgTpnt}{black}\colorlet{lgText}{black}
\else
  \colorlet{cPrime}{blue!12}\colorlet{dPrime}{black}\colorlet{tPrime}{black}
  \colorlet{cMu}{teal!15}\colorlet{dMu}{teal!60!black}\colorlet{tMu}{black}
  \colorlet{cLi}{yellow!25}\colorlet{dLi}{orange!70!black}\colorlet{tLi}{black}
  \colorlet{cGammaU}{blue!15}\colorlet{cGammaL}{red!15}\colorlet{dGamma}{black}\colorlet{tGamma}{black}
  \colorlet{cZetaU}{blue!15}\colorlet{cZetaL}{red!15}\colorlet{tZeta}{black}
  \colorlet{cXi}{cyan!15}\colorlet{dXi}{cyan!60!black}\colorlet{tXi}{black}
  \colorlet{cNT}{cyan!15}\colorlet{dNT}{cyan!60!black}\colorlet{tNT}{black}
  \colorlet{cProd}{cyan!15}\colorlet{dProd}{cyan!60!black}\colorlet{tProd}{black}
  \colorlet{cRH}{orange!15}\colorlet{dRH}{orange!70!black}\colorlet{tRH}{black}
  \colorlet{eArith}{blue!70!black}\colorlet{eElem}{violet!80!black}\colorlet{eDef}{black!80}
  \colorlet{ePNT}{brown!70!black}\colorlet{eAnal}{green!50!black}\colorlet{eExpl}{green!50!black}
  \colorlet{eExt}{red!70!black}\colorlet{eZero}{orange!80!black}\colorlet{eLate}{black!80}
  \colorlet{lgPrime}{blue!12}\colorlet{lgMu}{teal!15}\colorlet{lgLi}{yellow!25}
  \colorlet{lgXi}{cyan!15}\colorlet{lgRH}{orange!15}\colorlet{lgGU}{blue!15}\colorlet{lgGL}{red!15}
  \colorlet{lgEArith}{blue!70!black}\colorlet{lgEDef}{black!80}\colorlet{lgEElem}{violet!80!black}
  \colorlet{lgEZero}{orange!80!black}\colorlet{lgEAnal}{green!50!black}\colorlet{lgEPNT}{brown!70!black}\colorlet{lgEExt}{red!70!black}
  \colorlet{lgTprime}{black}\colorlet{lgTmu}{black}\colorlet{lgTli}{black}
  \colorlet{lgTxi}{black}\colorlet{lgTrh}{black}\colorlet{lgTgamma}{black}
  \colorlet{lgTarith}{black}\colorlet{lgTdef}{black}\colorlet{lgTelem}{black}
  \colorlet{lgTzero}{black}\colorlet{lgTanal}{black}\colorlet{lgTpnt}{black}\colorlet{lgText}{black}
\fi\fi
\makeatother

\begin{tikzpicture}[scale=1.0, transform shape,
  box/.style    = {draw, rounded corners=5pt, minimum width=2.8cm,
                   minimum height=0.9cm, align=center, font=\small, thick},
  zbox/.style   = {draw, rounded corners=5pt, minimum width=2.8cm,
                   minimum height=0.9cm, align=center, font=\small,
                   thick, fill=gray!18},
  splitbox/.style = {draw, rounded corners=5pt, minimum width=7cm,
                     minimum height=2.2cm, align=center, font=\small, thick,
                     rectangle split, rectangle split parts=2,
                     rectangle split part fill={red!8, white}},
  arr/.style    = {-{Stealth[length=6pt,width=5pt]}, thick},
  barr/.style   = {{Stealth[length=6pt,width=5pt]}-{Stealth[length=6pt,width=5pt]}, thick},
  % 线型：虚线仅用于 N(T)→RH、ξ↔RH；其余为实线
  proven/.style = {solid},
  cond/.style   = {dashed},
  lab/.style    = {font=\tiny, align=center, inner sep=3pt},
  rhbox/.style  = {draw=dRH, rounded corners=5pt, minimum width=8cm,
                   minimum height=1.0cm, align=center, font=\small,
                   fill=cRH, text=tRH, thick}
]

% 布局与 A-2 同步：左移 1.0cm、上移 2.0cm
% 左右略放宽，避免 Γ/ζ 自环旁公式标注被裁切
\path[use as bounding box]
  (-10.2, -26.2) rectangle (10.0, 7.4);

\begin{scope}[shift={(-1.0,2.0)}]

\node[font=\bfseries\Large, align=center] at (0, 4.35)
  {黎曼猜想知识谱系关系图};

% 主图单独上移 1cm；标题与图例位置不动
\begin{scope}[shift={(-2.75,-1.2)}]
%% ==================== 节点分层 ====================

% 第一层：素数计数函数 (y=0)
\node[box, fill=cPrime, draw=dPrime]    (J)    at (-4.2,  0.0) {$J(x)$};
\node[box, fill=cPrime, draw=dPrime]    (pi)   at ( 0.0,  0.0) {$\pi(x)$};
\node[box, fill=cPrime, draw=dPrime]    (th)   at ( 5.0,  0.0) {$\theta(x)$};
\node[box, fill=cPrime, draw=dPrime]    (ps)   at ( 9.7,  0.0) {$\psi(x)$};

% 第二层：积性函数 (y=-3.0)
% 节点定义修改
\node[box, fill=cMu, draw=dMu] (mu)  at (0.3, -3.0) {$\mu(n)$};
\node[box, fill=cMu, draw=dMu] (Lam) at (5.0, -3.0) {$\Lambda(n)$};

%Weierstrass乘积
\node[box, fill=cProd, draw=dProd, text=tProd]  (prod)   at ( 9.5,  -17.0) {\tiny\text{Hadamard 乘积}};
% M(x) 节点已移除：初等等价信息直接标注在 RH 框内

% 第三层：Li 与 Gamma (y=-6.0 至 -7.0)
\node[box, fill=cLi, draw=dLi] (Li) at (-2.0, -6.0) {$\Li(x)$};

% Gamma(s) 上下分色：上部标收敛域，下部标延拓
\node[splitbox, rectangle split draw splits=false, minimum width=5cm,
      minimum height=4.7cm,
      rectangle split part fill={cGammaU, cGammaL}] (gamma) at (7.0, -6.0)
    {
     \centering \color{tGamma}$\operatorname{Re}(s)>0$
     \nodepart{two}
     \raisebox{6pt}{\color{tGamma}$\mathbb{C}\setminus\{0,-1,-2,\dots\}$}
    };
% Gamma(s) 标注（固定在节点中线左侧）
\node[font=\small, left=-30pt, text=tGamma] at (gamma.west) {$\Gamma(s)$};

% 【修订3】ζ(s) 节点：下半部分补全函数方程左侧 ζ(s)=
\node[splitbox, rectangle split draw splits=false, minimum height=4cm,
      rectangle split part fill={cZetaU, cZetaL}] (zeta) at (0.3, -10.5)
    {\centering\small
     \color{tZeta}$\displaystyle
       \sum_{n=1}^{\infty}\frac{1}{n^s}
       =\prod_{p}\dfrac{1}{1-p^{-s}}
       \quad(\operatorname{Re}(s)>1)$%
     \vphantom{$\displaystyle\sum_{n=1}^{\infty}$}%
     \nodepart{two}
     % 修订：补全 ζ(s)= 使函数方程完整
     \textcolor{tZeta}{$\displaystyle\zeta(s)=
       2^s\pi^{s-1}
       \sin\!\left(\frac{\pi s}{2}\right)\Gamma(1-s)\,\zeta(1-s)$%
     \vphantom{$\displaystyle\sum_{n=1}^{\infty}\mathbb{C}\setminus\{1\}$}}
    };
% ζ(s) 标注


% ξ(s) 节点：下移至与 RH 同一水平，右侧对称于 M(x)
\node[box, fill=cXi, draw=dXi, text=tXi, minimum height=1.2cm]
  (xi) at (4.85, -17.0)
  {$\xi(s)$\\ \tiny 非平凡零点$\rho$与$\zeta(s)$相同\\
  \tiny $\xi(s)=\xi(1-s)$(函数方程推论)
  };


% 【修订8】N(T) 节点：置于 ζ 与 RH 正中间，构成竖直逻辑链
\node[box, fill=cNT, draw=dNT, text=tNT, minimum height=1.3cm]
  (NT) at (2.3, -14.0)
  {$N(T)$\\ \tiny $\zeta(s)$ 零点计数函数\\[1pt]\tiny（幅角原理）};

% ==================== 黎曼猜想节点 ====================
% RH 下移至 y=-20.5，为 N(T) 留出充分空间
\node[rhbox, minimum width=4.6cm, minimum height=2.2cm] (rh) at (-1.9, -17.0) {
    {\tiny\color{tRH} 黎曼猜想 (Riemann Hypothesis, 1859)}\\[3pt]
    {\tiny\color{tRH}$\zeta(s)$ 的所有非平凡零点 $\rho$ 满足
    $\operatorname{Re}(\rho)=\dfrac{1}{2}$}\\[4pt]
{\tiny\color{tRH} 初等等价（梅滕斯函数）：%
  $\displaystyle M(x)=\sum_{n\leq x}\mu(n)$}\\
{\tiny\color{tRH}
  $\mathrm{RH} \iff M(x) = O(x^{1/2+\varepsilon}),\ \forall \varepsilon>0$}
};


% ==================== 连线 ====================

% ----- 第一层内部 -----
\draw[barr, eArith] (J.east) -- (pi.west)
  node[lab, midway, above=14pt, xshift=-2pt, eArith]
    {$\displaystyle\pi(x)=\sum_{n=1}^{\infty}\frac{\mu(n)}{n}J(x^{1/n})$}
  node[lab, midway, below=14pt, xshift=-2pt, eArith]
    {$\displaystyle J(x)=\sum_{n=1}^{\infty}\frac{1}{n}\pi(x^{1/n})$};

\draw[barr, eElem] (pi.east) -- (th.west)
  node[lab, midway, above=14pt, eElem]
    {$\displaystyle\theta(x)=\pi(x)\log x-\int_2^x\frac{\pi(t)}{t}\,\mathrm{d}t$}
  node[lab, midway, above={14pt}, yshift={-50pt}, eElem]
    {\tiny $\displaystyle\pi(x) = \dfrac{\theta(x)}{\log x}+ \int_{2}^{x} \frac{\theta(t)}{t\,(\log t)^2} \,\mathrm{d}t$};

\draw[barr, eElem] (th.east) -- (ps.west)
  node[lab, midway, above=11pt, xshift={5pt}, eElem]
    {$\displaystyle\psi(x)=\sum_{k=1}^{\infty}\theta(x^{1/k})$}
  node[lab, midway, above=-31pt, xshift={5pt}, yshift={-8pt}, eElem]
    {\tiny $\displaystyle\theta(x)=\sum_{k=1}^{\infty}\mu(k)\psi(x^{1/k})$};


% ----- 第二层内部 -----
\draw[barr, eArith] (mu.east) -- (Lam.west)
  node[lab, midway, yshift=20, xshift=0.3cm, eArith]
    {$\displaystyle\Lambda(n)=-\sum_{d\mid n}\mu(d)\log d$
     \quad{\rm (M\"{o}bius 反演)}}
  node[lab, midway, yshift={-12pt-0.5cm}, xshift=0, eArith]
    {$\displaystyle\log n=\sum_{d\mid n}\Lambda(d)$};

% μ(n) → M(x) 的定义连线已省去（见节点内文字），避免与已有连线交叉

% Lambda -> psi
\draw[arr, eDef] (Lam.east) to[out=0, in=-90, looseness=1.0]
  ($(ps.south west)+(30pt,0)$)
  node[lab, pos=0.5, yshift=-90, xshift=250, eDef]
    {$\displaystyle\psi(x)=\sum_{n\le x}\Lambda(n)$};

% ----- 第三层：Li 和 Gamma -----

% 【修订5】Li(x) → π(x)：改为渐近等价语义（非因果推导）
\draw[arr, ePNT]
  ($(Li.north)!.75!(Li.north east)$)
  to[out=90, in=-135, looseness=1.8]
  ($(pi.south west)+(35pt,0)$)
  node[lab, pos=0.4, yshift=-45, xshift=45, ePNT, align=left]
    {\textbf{素数定理（PNT）}：$\pi(x)\sim\Li(x)$\\
     (若RH成立，$\pi(x)=\Li(x)+O(\sqrt{x}\log x)$)};

% Li -> J（显式公式主项；属 eExpl：第2/9章灰白，第15章全文总图着色）
\draw[arr, eExpl]
  (Li.north) to[out=135, in=-90, looseness=1.5]
  ($(J.south east)+(-22pt,0)$)
  node[lab, pos=0.5, yshift=-115, xshift=-65, eExpl, align=center]
    {$\Li(x)$ 是 $J(x)$ 显式\\[2pt]公式的主项};

% Gamma（Re>0）→ ζ：Euler 积分表示（该积分公式本身要求 Re(s)>1）
\draw[arr, eAnal]
  ([yshift=6pt] gamma.west)
  to[out=180, in=90, looseness=0.8]
  ($(zeta.north)+(30pt,0)$)
  node[lab, pos=0.5, yshift=-235, xshift=90,eAnal, align=center]
    {$\displaystyle\zeta(s)\Gamma(s)
      =\int_0^{\infty}\frac{x^{s-1}}{e^x-1}\,\mathrm{d}x$\\
     $(\operatorname{Re}(s)>1)$};

% μ(n) → ζ(s)
\draw[arr, eAnal]
  (mu.south) to[out=-90, in=90, looseness=1.2]
  ($(zeta.north)+({-20pt+0.3cm},0)$)
  node[lab, pos=0.5, yshift={-120pt-0.5cm-0.4cm}, xshift={45pt-5pt}, eAnal]
    {$\displaystyle\frac{1}{\zeta(s)}=\sum_{n=1}^{\infty}\frac{\mu(n)}{n^s}$\\
     $(\operatorname{Re}(s)>1)$};

% ----- J(x) 与 ψ(x) 的 Stieltjes 积分互逆关系 -----
% 【修订4】加注 dψ(t)/ln t 在 t=1 附近需 t>1 的说明
\draw[barr, eArith, line width=0.8pt, rounded corners=8pt]
  ($(J.north)+(0,0.0cm)$)
  -- ++(0,2.0cm)
  -- ++(13.5cm,0)
  -- ++(0,-2.0cm)
  node[lab, pos=0.35, above=6pt, yshift={10pt+0.5cm}, xshift=-200, eArith, align=center]
    {$\displaystyle\psi(x)=\int_{2^-}^x\log t\,\mathrm{d}J(t)$\quad(Stieltjes 积分，$J\to\psi$)}
  node[lab, pos=0.65, below=6pt, yshift={32pt+0.2cm}, xshift=-205, eArith, align=center]
    {$\displaystyle J(x)=\int_{2^-}^x\frac{\mathrm{d}\psi(t)}{\log t}$
     \quad(Stieltjes 反演，$\psi\to J$；积分自 $t>1$)};

% ----- 第四层：ζ 和 ξ 与第二层的连线 -----

% ζ → ψ：显式公式（Perron 积分 + 留数，单向；x≠p^k 时公式精确成立）
\draw[arr, eExpl, rounded corners=30pt]
  ($(zeta.east)+(0, -10pt)$) -| ($(ps.south)+(30pt,0)$)
  node[lab, pos=0.85, right=10pt, yshift=-170, xshift={-170pt-0.2cm},eExpl, align=left]
{$\psi(x)$ 显式公式（Perron 积分，$\zeta\to\psi$）\\
     $\displaystyle \psi(x) = \frac{1}{2\pi i}\int_{c-i\infty}^{c+i\infty}\!\left(-\frac{\zeta'(s)}{\zeta(s)}\right)\!\frac{x^s}{s}\,\mathrm{d}s$\\
     $\displaystyle = x-\sum_{\rho}\frac{x^{\rho}}{\rho}-\log(2\pi)-\tfrac{1}{2}\log(1-x^{-2})$\\
     \tiny $x>1,\;x\neq p^k$；$\rho$ 按 $|\mathrm{Im}(\rho)|$ 升序配对，条件收敛};

% ψ → ζ：Mellin 变换（单向）
\draw[arr, eExpl, rounded corners=30pt]
  ($(ps.south)+(40pt,0)$) |- ($(zeta.east)+(0, -20pt)$)
  node[lab, pos=0.8, yshift=-27, xshift=30,eExpl, align=center]
    {Mellin 变换（$\psi(x) \longleftrightarrow -\dfrac{\zeta'(s)}{\zeta(s)}$，即负对数导数）\\
     $\displaystyle -\frac{\zeta'(s)}{\zeta(s)} = s\int_{1}^{\infty}\frac{\psi(x)}{x^{s+1}}\,\mathrm{d}x$};
       
       
       
       
       
       
       
% Lambda -> zeta（对数微分 Dirichlet 级数）
\draw[arr, eAnal]
  ($(Lam.south west)+(6pt,0)$) to[out=-90, in=90, looseness=0.4]
  ($(zeta.north west)!0.50!(zeta.north east)$)
  node[lab, pos=0.5, yshift={-120pt-0.5cm-0.4cm+1cm}, xshift={145pt-5pt}, eAnal]
    {$\displaystyle-\frac{\zeta'(s)}{\zeta(s)}
      =\sum_{n=1}^{\infty}\frac{\Lambda(n)}{n^s}$};

% J(x) 显式公式（ζ → J）
\draw[arr, eExpl, rounded corners=50pt]
  ($(zeta.south west)+(0,6pt)$) -- ++(-60pt,0) -- ($(J.south west)+(5pt,0)$)
  node[lab, pos=0.5, yshift=-203, xshift={75pt-0.2cm}, eExpl, align=center]
{\textbf{Perron 公式} $\Rightarrow$ $J(x)$ 显式展开\\
     $\displaystyle
     \begin{aligned}
         &J(x)=\frac{1}{2\pi i}\int_{c-i\infty}^{c+i\infty}
           \frac{\log\zeta(s)}{s}\,x^s\,\mathrm{d}s\\
         &=\Li(x)-\sum_{\rho}\Li(x^{\rho})
           -\log 2+\int_x^{\infty}\frac{\mathrm{d}t}{t(t^2-1)\log t}
     \end{aligned}$\\
     \tiny $x>1,\;x\neq p^k$；$\rho$ 按 $|\mathrm{Im}(\rho)|$ 升序配对，条件收敛};

% ζ → J：Mellin 变换（ln ζ）
\draw[arr, eAnal, rounded corners=5pt]
  (zeta.north west) to[out=110, in=-90, looseness=1.5] (J.south)
  node[lab, pos=0.4, yshift=-205, xshift=-55, eAnal, align=center]
    {$\displaystyle\log\zeta(s)=s\int_1^{\infty}J(x)\,x^{-s-1}\,\mathrm{d}x$\\
     $(\operatorname{Re}(s)>1)$};

% ζ → ξ（定义）：ξ 已下移，连线沿右侧空白区域向右下延伸
\draw[arr, eLate, rounded corners=15pt]
  ($(zeta.south east)!0.15!(zeta.north east)$)
  -- ++(1.0, 0)
  -- ++(0, -3.8)
  -- ($(xi.north)+(0pt,0)$)
  node[lab, pos=0.25, right=4pt,yshift=30, xshift=0, eLate, align=center]
    {$\xi(s)=\dfrac{1}{2}s(s-1)\pi^{-s/2}\Gamma(s/2)\zeta(s)$  
};






% xi -> prod (Weierstrass product connection)
\draw[arr, eLate]
  %(xi.south) -- (prod.south)
  (xi.east) --(prod.west)
  node[midway, right=4pt, yshift=28, xshift=-45, eLate]
    {\tiny$\displaystyle\xi(s)=e^{A+Bs}\prod_{\rho}\left(1-\frac{s}{\rho}\right)e^{s/\rho}$};

% prod -> ps (zero sum in explicit formula)
%\draw[arr, violet!70!black, rounded corners=15pt]
  %(prod.east) -- ++(4.4, 0) |- (ps.south)
  %node[lab, pos=0.6, right=2pt, violet!70!black, align=left]
   % {对数微分与 Perron 积分\\
    % 给出显式公式中\\
    % 零点和 $\sum_{\rho}\frac{x_{\rho}}{\rho}$ 项};
%
% ζ 自身左侧弧线（解析延拓）：与表 tab:spectrum-stage8-guide 一致
\draw[arr, eExt, rounded corners=5pt]
  ($(zeta.north west)!0.15!(zeta.south west)$)
  to[out=180, in=180, looseness=1.6]
  ($(zeta.north west)!0.85!(zeta.south west)$);
% 标签用独立节点锚定在 ζ 西侧弧线旁（避免 path 末 node 漂移）
% 单行；中心在 zeta.west 右 2.5cm（相对初版再累计右移）
\node[lab, eExt, align=center, anchor=center]
  at ([xshift=2.5cm]zeta.west)
  {$\zeta$ 亚纯延拓至 $\mathbb{C}\setminus\{1\}$（Hankel / $\vartheta$ 两法）};

% ζ → N(T)：竖直向下（弧线）
\draw[arr, eZero]
  ($(zeta.south east)!0.10!(zeta.north east)$)
  to[out=0, in=90] ++(0.8, -0.8)
  to[out=-90, in=0] (NT.east)
  node[lab, pos=0.6, right=6pt, yshift=-357, xshift=5, eZero, align=center]
    {$\displaystyle N(T)=\frac{T}{2\pi}\log\frac{T}{2\pi e}+O(\log T)$\\
     \tiny 零点计数公式（von Mangoldt）};

% 【已删去】ζ → RH 蓝色直连线（设计意图改为ζ→N(T)→RH线性链）
%\draw[arr, eArith, rounded corners=10pt]
%  (zeta.south)
%  -- ++(0, -3.35)
%  -- ++(-4.1, 0)
%  -- ($(rh.north)!0.80!(rh.north west)$)
%  node[lab, pos=0.35, left=6pt, yshift=20pt,xshift=120,eArith, align=center]
%    {
%     $\zeta(s)$ 的非平凡零点全部位于直线 $\operatorname{Re}(s)=\dfrac12$ 上};
% RH -- P(s) 画线 + 文字
%\draw[arr, cyan!60!black]
 % ($(rh.east)!0.60!(rh.south east)$) -- (prod.west)    % 起点 -- 终点（必须有！）
  %node[midway, right=6pt, font=\tiny, cyan!60!black, align=center]
  %  {$\zeta(s) = \frac{e^{s(\log(2\pi) - 1)}}{2(s-1)} \prod_{n=1}^{\infty} \left(1 + \frac{s}{2n}\right)e^{-s/2n} P(s)$};

% ξ(s) ↔ RH、N(T) → RH：虚线标注（与 RH 相关的谱系连线）

% N(T) → RH：完成 ζ→N(T)→RH 线性逻辑链（虚线）
\draw[arr, eZero, cond, rounded corners=8pt]
  (NT.south west)
  to[out=-150, in=60, looseness=0.9]
  ($(rh.north)+(20pt, 0)$)
node[lab, pos=0.55, left=4pt, yshift=-430, xshift=-25, eZero, align=center]
    {$\mathrm{RH} \iff N_0(T)=N(T),\ \forall\,T>0$\\
   \tiny（$N_0(T)$：临界线 $\operatorname{Re}(s)=\dfrac{1}{2}$ 上虚部 $\leq T$ 的零点数）};

\draw[barr, eZero, cond]
  (xi.west) -- ($(rh.east)!0.0!(rh.south east)$)
  node[lab, midway, above=20pt,yshift=3pt,xshift=5, eZero, align=center]
    {$\mathrm{RH}\iff\xi\!\left(\dfrac{1}{2}+it\right)$\\
     的所有零点 $t$ 均为实数};

% M(x)→RH 箭头已删去：等价信息直接写在 RH 框内

% Gamma 自身解析延拓弧线：与表 tab:spectrum-stage8-guide 一致
\draw[arr, eExt, rounded corners=1.5pt]
  ($(gamma.north east)+(0,-5pt)$)
  to[out=0, in=0, looseness=1.6]
  ($(gamma.south east)+(0,5pt)$);
% 标签用独立节点锚定在 Γ 东侧弧线旁（仅递推式；延拓域已在框下半）
% 相对初版 (gamma.east+14pt)：净向左 3.7cm，上移 2pt
\node[lab, eExt, align=center, anchor=west]
  at ($(gamma.east)+({14pt-3.7cm},2pt)$)
  {递推延拓 $\Gamma(s)=\Gamma(s+1)/s$};

% Gamma（延拓后）→ ζ：Hankel 围道
\draw[arr, eExt]
  ($(gamma.south west)+(40pt,0)$)
  to[out=180, in=0, looseness=1.0]
  ($(zeta.south east)!0.45!(zeta.north east)$)
  node[lab, pos=0.5, yshift=-210, xshift=215, eExt, align=center]
    {$\displaystyle\zeta(s) = -\frac{\Gamma(1-s)}{2\pi i}\oint_{\mathrm{H}}\frac{(-z)^{s-1}}{e^z-1}\,\mathrm{d}z$\\
     解析延拓到 $\mathbb{C}\setminus\{1\}$，Hankel 围道消去分支切割};

\end{scope}% 主图

% ==================== 图例（取自 A-2 详细文案 + 分阶段灰白色） ====================
\node[draw=black!40, fill=gray!8, rounded corners=8pt, thick, align=center, inner sep=10pt, font=\tiny]
  (legend) at (0, -23.5) {
  \begin{tabular}{@{}ll@{\hspace{1.0cm}}ll@{}}
    \multicolumn{2}{l}{\scriptsize\textbf{节点颜色分类图例}} &
    \multicolumn{2}{l}{\scriptsize\textbf{知识谱系连线语义图例}} \\[6pt]
    \tikz[baseline=-0.5ex]\fill[lgPrime, draw=black!70, thick, rounded corners=1.5pt] (0,-0.12) rectangle (0.55,0.12); &
    \textcolor{lgTprime}{\textbf{素数计数函数}（实变，如 $\pi(x)$）} &
    \textcolor{lgEArith}{\rule[0.5ex]{16pt}{1.4pt}} &
    \textcolor{lgTarith}{\textbf{算术互逆}（M\"{o}bius / Stieltjes 反演）} \\[3pt]
    \tikz[baseline=-0.5ex]\fill[lgMu, draw=black!70, thick, rounded corners=1.5pt] (0,-0.12) rectangle (0.55,0.12); &
    \textcolor{lgTmu}{\textbf{算术积性函数}（离散，如 $\mu(n)$）} &
    \textcolor{lgEDef}{\rule[0.5ex]{16pt}{1.4pt}} &
    \textcolor{lgTdef}{\textbf{对象定义}（构造定义与主项提取）} \\[3pt]
    \tikz[baseline=-0.5ex]\fill[lgLi, draw=black!70, thick, rounded corners=1.5pt] (0,-0.12) rectangle (0.55,0.12); &
    \textcolor{lgTli}{\textbf{渐近逼近主项}（如 $\Li(x)$）} &
    \textcolor{lgEElem}{\rule[0.5ex]{16pt}{1.4pt}} &
    \textcolor{lgTelem}{\textbf{初等计数变换}（数论函数基本恒等式）} \\[3pt]
    \tikz[baseline=-0.5ex]\fill[lgXi, draw=black!55, thick, rounded corners=1.5pt] (0,-0.12) rectangle (0.55,0.12); &
    \textcolor{lgTxi}{\textbf{辅助复变函数}（如 $\xi(s)$）} &
    \textcolor{lgEZero}{\rule[0.5ex]{16pt}{1.4pt}} &
    \textcolor{lgTzero}{\textbf{零点结构与猜想}（分布特征及等价命题）} \\[3pt]
    \tikz[baseline=-0.5ex]\fill[lgRH, draw=black!55, thick, rounded corners=1.5pt] (0,-0.12) rectangle (0.55,0.12); &
    \textcolor{lgTrh}{\textbf{核心猜想与目标}（如 RH）} &
    \textcolor{lgEAnal}{\rule[0.5ex]{16pt}{1.4pt}} &
    \textcolor{lgTanal}{\textbf{分析桥梁}（积分变换与 Dirichlet 级数）} \\[3pt]
    \tikz[baseline=-0.5ex]{\clip[rounded corners=1.5pt] (0,-0.12) rectangle (0.55,0.12); \fill[lgGL] (0,-0.12) rectangle (0.55,0); \fill[lgGU] (0,0) rectangle (0.55,0.12); \draw[black!55, thick, rounded corners=1.5pt] (0,-0.12) rectangle (0.55,0.12); \draw[black!55, thick] (0,0) -- (0.55,0);} &
    \textcolor{lgTgamma}{\textbf{解析延拓复变}（初始域/延拓域）} &
    \textcolor{lgEPNT}{\rule[0.5ex]{16pt}{1.4pt}} &
    \textcolor{lgTpnt}{\textbf{渐近定理}（PNT 等确立的渐近推论）} \\[3pt]
    & & \textcolor{lgEExt}{\rule[0.5ex]{16pt}{1.4pt}} &
    \textcolor{lgText}{\textbf{解析延拓}（复围道积分与全纯延拓）} \\[3pt]
    & & \tikz[baseline=-0.5ex]{\draw[lgEZero, line width=1.2pt] (0.00,0) -- (0.12,0); \draw[lgEZero, line width=1.2pt] (0.17,0) -- (0.29,0); \draw[lgEZero, line width=1.2pt] (0.34,0) -- (0.46,0);} &
    \textbf{未证}
  \end{tabular}
};

\end{scope}% 位移：左 1.0cm、上 2.0cm

\end{tikzpicture}
%
}
\end{figure}
\clearpage
\noindent
\captionof{figure}{黎曼猜想知识谱系关系图（第三次着色：全文）。
与图~\ref{fig:spectrum-stage2}（第一次）、图~\ref{fig:spectrum-stage8}（第二次）同一底图；
彩色表示截至第~\ref{ch:riemann_explicit_formula}~章已建立的主要对象与关系。
相对第二次进度新增的彩色节点/连线与正文对照见表~\ref{tab:spectrum-stage16-guide}。}
\label{fig:function-relations}

\paragraph{图~\ref{fig:function-relations} 新增彩色元素与正文对照}
下列各项将图上\textbf{相对第二次进度图新着色}的节点与连线引向相应定义与证明；
第一次、第二次进度分别见表~\ref{tab:spectrum-stage2-guide}、
表~\ref{tab:spectrum-stage8-guide}，此处不重复。

\begin{table}[htbp]
\centering
\caption{图~\ref{fig:function-relations} 新增彩色图元与正文对照}
\label{tab:spectrum-stage16-guide}
\small
\renewcommand{\arraystretch}{1.35}
\begin{tabular}{@{}p{4.9cm}p{8.3cm}@{}}
\toprule
\textbf{图上元素} & \textbf{正文（定义 / 公式 / 证明）} \\
\midrule
节点 $\xi(s)$ &
  第~\ref{ch:xi_function}~章第~\ref{sec:xi-definition}~节：
  定义 $\xi(s)=\frac12 s(s-1)\pi^{-s/2}\Gamma(s/2)\zeta(s)$，
  对称形式 $\xi(s)=\xi(1-s)$ \\
连线 $\zeta\to\xi$（定义） &
  同上；由第~\ref{ch:riemann_zeta_continuation}~章函数方程消去极点与附加因子 \\
节点 Hadamard 乘积 &
  第~\ref{ch:xi_function}~章第~\ref{sec:xi-hadamard-product}~节：
  $\xi(s)=e^{A+Bs}\prod_{\rho}(1-s/\rho)\,e^{s/\rho}$ \\
连线 $\xi\to$ Hadamard 乘积 &
  同上：一级整函数的 Weierstrass--Hadamard 因子分解 \\
节点 $N(T)$ &
  第~\ref{ch:zero_distribution}~章第~\ref{sec:NT-RVM}~节：
  零点计数函数与 Riemann--von Mangoldt 公式 \\
连线 $\zeta\to N(T)$
  （$N(T)=\frac{T}{2\pi}\log\frac{T}{2\pi e}+O(\log T)$） &
  第~\ref{sec:NT-RVM}~节：幅角原理与零点阶梯计数 \\
RH 框（非平凡零点在 $\operatorname{Re}s=\frac12$） &
  第~\ref{ch:zero_distribution}~章第~\ref{sec:critical_line_rh}~节：
  猜想陈述；临界带限制见第~\ref{sec:zero-distribution-strip}~节 \\
连线 $N(T)\to\mathrm{RH}$
  （$N_0(T)=N(T)$ 等价） &
  第~\ref{sec:critical_line_rh}~节：
  临界线上零点计数 $N_0(T)$ 与 $N(T)$ 全等 \\
连线 $\xi\leftrightarrow\mathrm{RH}$
  （$\xi(\frac12+it)$ 的零点 $t$ 皆实） &
  第~\ref{sec:critical_line_rh}~、\ref{sec:two-languages}~节：
  $\xi$ 语言下的 RH 等价表述 \\
RH 框内 $M(x)$ 初等等价 &
  第~\ref{ch:riemann_explicit_formula}~章第~\ref{sec:mertens-rh}~节：
  $\mathrm{RH}\iff M(x)=O(x^{1/2+\varepsilon})$ \\
\midrule
连线 $\zeta\to\psi$（Perron 显式公式） &
  第~\ref{ch:riemann_explicit_formula}~章第~\ref{sec:explicit-derivation}~节式~\eqref{eq:explicit_psi}；
  Perron 核见第~\ref{ch:integral_transforms}~章第~\ref{sec:perron}~节 \\
连线 $\psi\to\zeta$
  （Mellin：$-\zeta'/\zeta=s\int\psi(x)x^{-s-1}\,\mathrm{d}x$） &
  第~\ref{ch:riemann_explicit_formula}~章式~\eqref{eq:psi_mellin}；
  一般形式见第~\ref{ch:integral_transforms}~章 \\
连线 $\zeta\to J$
  （Perron $\Rightarrow$ $J$ 显式展开） &
  第~\ref{ch:riemann_explicit_formula}~章第~\ref{sec:integration-by-parts}~节式~\eqref{eq:explicit_J_final} \\
\bottomrule
\end{tabular}
\end{table}

% ============================================================
\section{读图要点}
\label{sec:spectrum_reading}
% ============================================================

读图时，箭头应读作\textbf{已证公式的索引}，而非替代各章证明。
图上主线可概括为：
算术计数（$\pi,J,\theta,\psi,\mu,\Lambda$）$\to$ 分析对象（$\zeta,\Gamma$）$\to$
整函数与零点（$\xi$、Hadamard、$N(T)$）$\to$ 显式公式（$\psi,J$ 与零点级数）$\to$ 黎曼猜想（RH）；
RH 是横跨“零点—误差”的最强假设，而非谱系终点上的已证推论。

此后各章转向振荡的 Fourier 视角、数值与几何示意、1859 年原稿解读及当代前沿，
\textbf{不再}向该底图追加主体理论节点。
（书名何以将「素数分布」与「黎曼猜想」并列，见书末后记。）

% ============================================================
\section*{本章小结与后续展望}
\addcontentsline{toc}{section}{本章小结与后续展望}
% ============================================================

本章给出与第~\ref{ch:elementary_number_theoretic_functions}~、
\ref{ch:riemann_zeta_continuation}~章同一底图的\textbf{全文着色}总图~\ref{fig:function-relations}，
并附新增彩色图元与正文对照表。
主体理论地图至此齐备；下一章起讨论振荡的 Fourier 重组与逼近比较。
